% =====================================================================
%  H. Brøchner
%  Et Svar til Professor R. Nielsen.
%  Kjøbenhavn: P. G. Philipsens Forlag, 1868.   31 pp.
%  Printer: Bianco Lunos Bogtrykkeri ved F. S. Muhle.
%
%  DIPLOMATIC TRANSCRIPTION (Danish).
%  Copy digitised: BRITISH LIBRARY copy, scanned by Google in 2013,
%  Google Books id hnpZAAAAcAAJ. The Royal Danish Library has NOT
%  digitised this pamphlet: KB records only the Aarhus copy (s Teol
%  Brøchner H., læsesalsbrug) and a 2. oplag bound into a Nielsen
%  samlingsbind at KU Teologi (V g 2 Nie). Neither was consulted; if the
%  two printings ever turn out to differ, this edition represents the
%  British Library copy of the first.
%
%  The last round of the Brøchner--R. Nielsen exchange. It answers
%  Nielsen's »Hr. Professor Brøchner's philosophiske Kritik gjennemseet«
%  (1868), which was itself a reply to Brøchner's »Problemet om Tro og
%  Viden« (1868) -- both in this repository.
% =====================================================================
%
%  ---------------------------------------------------------------
%  1. PAGE MAP  (verified by eye, recon pass, 4 September 2026)
%  ---------------------------------------------------------------
%  The scan is 42 PDF pages. Each page is a single 600 ppi 1-bit JBIG2
%  image (about 3016x4842 px) plus a 1034x204 Google watermark strip at
%  the foot. The map is UNIFORM over the whole book:
%
%      printed p. N   =   PDF page N + 4        (printed 1--31 = PDF 5--35)
%
%  Verified by reading the printed folio off the image at printed pp. 4,
%  6, 10, 11, 20, 24, 26, 29 and 31, and corroborated mechanically by an
%  ABBYY folio sweep of all 29 body leaves, which contradicts +4 nowhere.
%  Three body leaves carry no folio by design -- PDF 5 (p. 1, title page),
%  PDF 6 (p. 2, printer's imprint) and PDF 7 (p. 3, the opening page of
%  the text) -- and each is fixed from below by PDF 8 = printed p. 4.
%
%  The rest of the file, none of it part of the pamphlet:
%      PDF  1   Google digitisation notice
%      PDF  2   blank
%      PDF  3   printed WRAPPER (yellow paper), the title repeated inside
%               a double frame, from the same setting as the title page
%               but with a single rule where the title page has a double
%      PDF  4   wrapper verso: P. G. Philipsens Forlag advertisements
%      PDF 36   British Library date stamp „10 JU 69“
%      PDF 37   the RETTELSE leaf  (transcribed -- see §5)
%      PDF 38   ditto date stamp
%      PDF 39,40,42  P. G. Philipsens Forlag advertisements
%      PDF 41   blank
%  Also not text: the British Museum oval accession stamp on PDF 6, and
%  the modern pencil marks („3.“ on the wrapper, „7/“ on the title page)
%  which belong to this copy's later history.
%
%  ---------------------------------------------------------------
%  2. WHAT IS TRANSCRIBED, AND HOW
%  ---------------------------------------------------------------
%  Diplomatic. 19th-century orthography exactly as printed: Christendommen,
%  Videnskaben, Theologie, høiere, aa, capitalised nouns. Printer's errors
%  are transcribed AS PRINTED, logged in a % comment at the site, and never
%  silently corrected -- including the one erratum the author himself lists
%  (see §5). Quote-balance counts therefore drift off zero on purpose.
%
%  The book is FRAKTUR throughout -- body, notes, title page and wrapper --
%  with the long s distinguished. Long s is normalised to s, the house
%  convention of the sibling volumes. There is no bold inside a paragraph
%  and no italic Fraktur anywhere: display lines and the running-head
%  folios are the only bold in the book.
%
%  STRUCTURE. There is none to speak of: no chapters, no sections, no
%  headings, no contents page, no block quotations, no lists, no tables.
%  The pamphlet is 29 pages of continuous polemical prose, printed 3--31.
%  A mechanical sweep of all 29 body pages for short centred lines found
%  nothing but the running head, which is the folio alone (absent by design
%  on pp. 1--3). Signature numerals „2“ (p. 17) and „2*“ (p. 19) are not
%  transcribed. Three centred rules ARE text and ARE transcribed -- pp. 26
%  and 29 (short hairlines, about 1.9 cm) and p. 31 (an ornamental
%  tailpiece, about 3.4 cm, closing the book) -- each with its printed form
%  recorded in a % comment at the site. The footnote separator rule is not
%  transcribed; it is LaTeX's own job.
%
%  QUOTATION MARKS. „…“, low open / high close, throughout -- the same
%  practice as the 1868 and 1869 sibling volumes. Two complications, both
%  kept as printed:
%    (a) DOUBLED „„…““ for a quotation inside a quotation, on p. 20, twice
%        (the Madvig grammar quotations). Single-level nesting also occurs,
%        e.g. p. 19 „Nielsens Exempler med de „mathematiske Venner“ …“.
%    (b) ANTIQUA GUILLEMETS «…» round antiqua foreign matter, following the
%        FOUNT and not the sense: «lazzi» (p. 19) and «practica … arithmetica»
%        (p. 21) carry them, while most of the antiqua Latin carries none --
%        calendæ græcæ (p. 3), Nomina sunt odiosa! (p. 5), instar omnium
%        (pp. 10, 18), lucus a non lucendo (p. 12), petitio principii
%        (p. 22), difficile est satiram non scribere (p. 27).
%  Em-dash: ---. One long sort (about 1 em, 84 px at 600 dpi against a body
%  line pitch of 101 px) serves both the parenthetical dash and the dash
%  between numerals (141---148, 1860---61).
%
%  EMPHASIS. Letterspacing (Sperrsatz) is the book's only emphatic device
%  inside a paragraph, and it is used very heavily -- whole phrases, whole
%  sentences, and once a seven-line quotation (the Erasmus Montanus motto,
%  p. 26). All of it is rendered \emph{}.
%
%  ANTIQUA WITHIN THE FRAKTUR: Latin and Greek only, never a proper name.
%  Not one proper name was found set in antiqua anywhere in the book;
%  names are letterspaced Fraktur instead. The printed Latin is UPRIGHT
%  antiqua and is rendered \textit{} here, for consistency with the sibling
%  volumes -- an editorial choice, recorded as such. The antiqua may itself
%  be letterspaced: spaced antiqua is \emph{\textit{…}}, tight antiqua
%  \textit{…}, and the two occur inside one phrase on p. 12, which prints
%  „lucus a n o n lucendo“ with only „non“ spaced.
%
%  GREEK: two words, both antiqua italic with accents -- ἐνέργεια (p. 14)
%  and ἄπειρον (p. 22). The Fraktur OCR model reads no Greek at all; the
%  ABBYY layer is the only witness for it and its breathings and accents
%  are not to be trusted, so both were read on the image at 600 dpi and
%  under magnification (the smooth breathing and acute of ἐνέργεια, the
%  combined breathing and acute of ἄ). Following the sibling volume
%  `det-religioese', the Greek is TYPED DIRECTLY and left to `textalpha'
%  rather than wrapped in \textit{}: the Greek fount is already distinct
%  from the Fraktur, so the italic marker would add nothing and would
%  render inconsistently across the two volumes.
%
%  The old id-est mark ɔ: (a reversed c followed by a colon) occurs; it is
%  typed as U+0254 + colon and rendered \reflectbox{c}.
%
%  The section sign is used inconsistently and is kept so: p. 22 has „§ 11“,
%  and p. 26 has „§ 35“ and „§ 36“ alongside „§. 11“ WITH A FULL STOP, in
%  the same paragraph.
%
%  NORMALISATION (the one departure from strict diplomacy, besides the
%  sort-level policy of TRANSCRIPTION-PLAYBOOK §2): Fraktur has a single
%  sort for I and J, so the print cannot distinguish them and the choice is
%  the editor's. It is resolved by the word -- I for I-words (Indsigelse,
%  Intet, „I den Efterskrift“), J for J-words (Jeg, Jfr.) -- following the
%  house convention of `problemet-tro-viden'.
%
%  ---------------------------------------------------------------
%  3. FOOTNOTES
%  ---------------------------------------------------------------
%  The book prints *) **) ***), restarting at *) on EVERY printed page, so
%  \setcounter{footnote}{0} stands at the page marker of every note-bearing
%  page. (Not footmisc[perpage], which resets on the TYPESET page.) Three
%  levels are actually used, on p. 6.
%
%  Fifteen notes stand on eleven pages: 3, 4, 5, 6, 9, 10, 15, 17, 18, 24
%  and 28. TWO OF THEM RUN OVER A PAGE BOUNDARY IN THE PRINT -- p. 4's **)
%  continues on p. 5, and p. 10's *) continues on p. 11, where it occupies
%  most of the leaf (p. 10's note fills some 33 of that page's 40 lines).
%  Pages 5 and 11 accordingly OPEN with an unmarked continuation in the
%  note area. Each such note is set here IN FULL at its call, on p. 4 and
%  p. 10 respectively, with a % comment at the p. 5 / p. 11 page break
%  recording that the note continues there in the print.
%
%  ---------------------------------------------------------------
%  4. THE SCAN, AND WHAT MAY NOT BE CLAIMED FROM IT
%  ---------------------------------------------------------------
%  This is a Google Books JBIG2 scan with symbol-matching compression --
%  the same encoder as the Spinoza volume in this repository, where 57--59 %
%  of the glyph components on a sample page proved pixel-identical to
%  another instance of the same letter. Such a file RECONSTRUCTS its letters
%  from a dictionary rather than reproducing each impression, so NO
%  SORT-LEVEL CLAIM MAY BE MADE FROM IT: no wrong-sort inventory, no
%  turned-sort or broken-sort log, no argument about what the compositor's
%  case held. An apparent turned n may be the encoder's substitution. This
%  edition therefore makes none. Structural defects -- dropped words, wrong
%  punctuation, transposed lines -- remain fair game and are logged where
%  found. See texts/brochner/spinoza/JBIG2-TEST.md for the experiment.
%
%  ---------------------------------------------------------------
%  5. ERRATA AND DEFECTS
%  ---------------------------------------------------------------
%  THE AUTHOR'S OWN RETTELSE, printed on a leaf at the back (PDF 37):
%
%      R e t t e l s e .
%      Side 11, Lin. 2: videnskabelige, læs: subjective.
%
%  It is right: p. 11 line 2 reads „Almindeliges videnskabelige Gyldighed
%  med den objective,“ where the sense plainly wants „subjective … med den
%  objective“. Per the repo's diplomatic stance the body KEEPS
%  „videnskabelige“ as printed, with a % comment at the site, and the
%  Rettelse is set at the end of this edition as the author printed it.
%  It is not silently applied.
%
%  OTHER PRINTER'S DEFECTS, all transcribed AS PRINTED with a % comment at
%  the site. Complete list for the finished edition:
%
%      p. 16   „poa“ for „paa“
%      p. 20   „Adjertiv“ for „Adjectiv“, inside the quotation from Nielsen
%      p. 20   the OUTER quotation, opened at „Istedenfor (S. 50. Anm.), is
%              never closed. This is the whole of the edition's quote
%              imbalance -- see below.
%      p. 21   „geometica“ for „geometrica“ (in the Holberg quotation)
%      p. 21   the opening „ set tight against „samledes“, i.e. „samledes„ paa“
%      p. 27   „Prof. N.'s vil kun“, where the sense wants „Prof. N. vil kun“
%      p. 24   note **) closes without a full stop
%
%  Both of the two readings that look like wrong sorts -- „poa“ and
%  „Adjertiv“ -- were calibrated at 600 dpi and 8x against certain
%  instances of the competing letter ON THE SAME PAGE before being logged,
%  per TRANSCRIPTION-PLAYBOOK §2. They are recorded as IMAGE READINGS, not
%  as claims about the compositor's case: see §4 above on why this file
%  cannot support a sort-level claim.
%
%  INK SPECKS AND BLOTS, examined and NOT transcribed, because each is far
%  below sort size and stands in a word space: p. 12 (a solid round blot in
%  „Væsen ud|trykker“), p. 13, p. 14 (twice), pp. 23, 24 (three), p. 25.
%  They are logged at the site because the Fraktur OCR reads several of
%  them as full stops.
%
%  QUOTE BALANCE. The edition ends at „=130, “=129, i.e. +1, and «=2, »=2.
%  THE +1 IS DELIBERATE and is the unclosed outer quotation on p. 20. Do
%  not "fix" it. „check.py“ prints these figures every run; a change in the
%  delta means something moved, and the running total is kept in
%  RESUME-NOTES.md.
% =====================================================================

\documentclass[12pt,a4paper]{book}
\usepackage[utf8]{inputenc}
\usepackage[T1]{fontenc}
\usepackage{libertinus}
\usepackage{libertinust1math}
\usepackage{textalpha}   % Greek (ἐνέργεια p. 14, ἄπειρον p. 22)
\usepackage{graphicx}    % for \reflectbox (the old Danish »ɔ:« id-est mark)
\DeclareUnicodeCharacter{0254}{\reflectbox{c}}  % ɔ (U+0254) = reversed c
\usepackage[danish]{babel}
\usepackage[protrusion=true,expansion=true]{microtype}
\usepackage{setspace}
\usepackage[margin=1.2in]{geometry}
\usepackage{enumitem}
\usepackage{fancyhdr}
\setlength{\headheight}{28pt}
\addtolength{\topmargin}{-13.5pt}
\usepackage{hyperref}

\hypersetup{
  pdftitle={Et Svar til Professor R. Nielsen},
  pdfauthor={H. Brøchner},
  hidelinks
}

% The book prints *) **) ***), restarting on every printed page.
% All three levels are used, on p. 6.
\renewcommand{\thefootnote}{%
  \ifcase\value{footnote}\or *)\or **)\or ***)\else \arabic{footnote})\fi}

\pagestyle{fancy}
\fancyhf{}
\fancyhead[LE,RO]{\thepage}
\fancyhead[RE]{\textit{Et Svar til Professor R. Nielsen}}
\fancyhead[LO]{\textit{H. Brøchner}}

\setstretch{1.3}

\title{\textbf{Et Svar\\[2pt] til\\[2pt] Professor R. Nielsen}}
\author{Dr.\ H.\ Brøchner\\[4pt]
  \small Professor i Philosophie}
\date{Kjøbenhavn: P.\ G.\ Philipsens Forlag, 1868}

\begin{document}
\maketitle
\thispagestyle{empty}
\newpage

% =====================================================================
%  TEXT.  Printed pp. 3--31, continuous prose, no divisions.
%  Printed p. 1 is the title page and p. 2 the printer's imprint
%  (Bianco Lunos Bogtrykkeri ved F. S. Muhle); neither is transcribed as
%  body text -- the title page is represented by \maketitle above.
% =====================================================================

% --- p. 3 ---
% Printed p. 3 opens with an oversized ornamental Fraktur initial, about twice
% body height and standing on the first baseline (not a sunk drop-cap). The
% Fraktur sort serves both I and J; the word is „I den Efterskrift“, so it is
% resolved as I and set as ordinary type, per the house convention.
\setcounter{footnote}{0}
I den Efterskrift, som Hr. Professor Nielsen har føiet til sine „populære“
Forelæsninger fra Christiania, slutter han med en Erklæring, affattet i myndig
Form og luxuriøst udstyret med Gadesprogets Bijouterier. Professor Nielsen vil
ikke svare, om ogsaa en femte og en sjette Angriber skulde slutte sig til de
fire, som der ikke skal være nogen Mening i at svare, og af hvilke dog de to, de
to ordinerede Mænd, maaskee ifølge Professorens gamle, og bestandig
tilbagevendende Kjærlighed til Theologien, --- vare blevne svarede.
% „majestætisk“: the æ is plain at 400 dpi. The ABBYY layer reads „majestetisk“
% and is wrong here.
Professor Nielsen drager sig med hiin Erklæring majestætisk tilbage, og han vil
nu, thronende paa sin tvetoppede Logiks Olymp, og indhyllet i sin Troestheories
Taageskyer, idethøieste udslynge, ikke just Skysamlerens Lyn, men den Art
Projectiler, der fordrive Fjenden, skjøndt de hverken kunne saare eller skade.

Saa kom den „femte Angriber“ med sit Skrift, og Professor Nielsen skyndte sig,
uden at vente paa Skriftets Afslutning, der dog ikke, som Prof. N.'s
Religionsphilosophie, først skulde oppebie Grundideernes Logiks Afslutning om
% The fraction is printed as a superscript 1, a fraction bar and a subscript 2;
% the ABBYY layer mangles it to „1142“.
114½ Aar\footnote{% printed *)
  Beregningen er gjort i Dr. Heegaards Skrift: Prof. R. N.'s Lære om Tro og
  Viden, S. 94 f.}
% „Dr.“ is ANTIQUA in the print here (and again on p. 11), as the arabic
% figures are and as the roman numerals „IX og X“ on p. 8 are. It is an
% abbreviation, not quoted foreign matter, so it is left plain rather than
% \textit{}-ed, following the treatment of the antiqua numerals in
% brochner/det-religioese.
og saa udkomme til de nærmeste \textit{calendæ græcæ}, men som skulde gives i
dette Aar; og han skrev en Piece, der bærer Titelen: Hr. Professor Brøchners
philosophiske Kritik, \emph{gjennemseet} af R. Nielsen.
% EMPHASIS, checked at 400 dpi and again at 800: within the cited title only
% „g j e n n e m ſ e e t“ is letterspaced. „Hr. Professor Brøchners
% philosophiske Kritik,“ on the line above and „af R. Nielsen.“ after it are
% both set tight. (RECON.md §5 reads the whole final line as spaced; the image
% does not bear that out.)

% --- p. 4 ---
\setcounter{footnote}{0}
Er denne Piece nu et Svar eller blot et Projectil, eller er den en høiere
dialektisk Enhed af Svar og Ikke-Svar, af Svar og Projectil? Det Efterfølgende
vil give Læseren de Præmisser, af hvilke han saa selv maa uddrage Conclusionen.

I min Kritik af Professor Nielsens Theorie havde jeg godtgjort, at den
Bestemmelse, der var Theoriens Hovedhjørnesteen: den \emph{absolute
Uensartethed}, som Prof. Nielsen saa stærkt accentuerede, og som Ecchoerne
gjentoge i Chor, udtrykte et Selvmodsigende og Umuligt.
% „den“ is set tight before the spaced „a b ſ o l u t e  U e n s a r t e t h e d“,
% here and in note *) below.

Jeg havde viist, hvorledes dette Begreb, som Prof. Nielsen ikke kunde fastholde
og ikke vilde opgive, bragte Forvirring ind overalt, og af sig fødte en Yngel af
Modsigelser, som jeg opregnede i mit Skrift S. 213 f.\footnote{% printed *)
  Prof. N. indrømmer (Piecen S. 3 nederst), at jeg har paaviist
  \emph{virkelige} Modsigelser, men han undgaaer at sige, hvilke. Skulde det
  maaskee blandt Andet være den \emph{absolute Uensartethed}?}.

Jeg havde viist, hvorledes den Nielsenske Theorie, medens den undergraver Videns
Grund, tilintetgjør det Ethiske og forvansker det Religiøse.

Jeg havde viist, hvorledes denne Theorie hos Prof. N. er et tomt
Tankeexperiment, uden Rod i en Personlighed, uden Sandhed i et Livs Virkelighed.

Jeg havde, med Hensyn til det videnskabelige Værk, til hvilket Prof. N. viser
tilbage som Grundlaget for sin Troestheorie, godtgjort dets Dualisme, og
paaviist en gjennemgaaende Forvirring i Grundbegreberne, en Mangel paa
videnskabelig Form. Jeg havde godtgjort Umuligheden af den projecterede
„omvendte“ Videnskab: Religionsphilosophien (i Nielsensk Betydning), den, ved
hvilken en fjern Fremtid skulde lyksaliggjøres.

Paa alt dette har Prof. Nielsen \emph{Intet} kunnet svare\footnote{% printed **)
  % THIS NOTE RUNS OVER THE PAGE BREAK IN THE PRINT: it begins in the note area
  % of p. 4 and its last four lines stand, unmarked, at the head of the note
  % area of p. 5, above that page's own *). It is set here IN FULL at its call,
  % per RECON.md §4; see the % comment at the p. 5 page marker below.
  Det er ligesom om en drillende Dæmon, ubevidst for Prof. N., har aflokket ham
  denne Tilstaaelse. I de allerførste Linier af sit Skrift indrømmer Prof. N.,
  at jeg har \emph{gjendrevet} adskillige andre Theorier. De eneste Theorier,
  jeg har behandlet foruden Nielsens, ere Martensens og Kierkegaards. Men nu har
  jeg gjort det tydeligt, at Nielsens Theorie staaer i et saadant Forhold til
  Kierkegaards, at Gjendrivelsen af den sidste ogsaa er Gjendrivelsen af den
  første, medens ganske vist Forholdet ikke kan vendes om. Hiin Indrømmelse
  gjorde det altsaa overflødigt for Prof. N. at skrive 52 Sider, han kunde nøies
  med 3 Linier.}:
% --- p. 5 ---
% The note area of p. 5 OPENS with the unmarked last four lines of p. 4's **)
% („Gjendrivelsen af den sidste … kunde nøies med 3 Linier.“); they are set
% above, at the call on p. 4, and are NOT repeated here. The page's own note,
% *) „Nomina sunt odiosa!“, follows them on the leaf and is set below.
\setcounter{footnote}{0}%
det staaer, efterat hans Piece er udkommen, som før, ugjendrevet. Og dog, naar
blot Læren om den \emph{absolute Uensartethed} maa opgives af Prof. N., saa er
Slaget tabt for ham; thi denne \emph{absolute} Uensartethed, „paa Grund af
hvilken“ Forsoningen af Tro og Viden i samme Bevidsthed skal blive mulig, bærer,
som tidt nok indrømmet, den \emph{hele Lære}. \emph{Men selve Ordet er undgaaet
i Prof. N.'s Skrift.}
% Letterspacing here is fine-grained and was read at 400 dpi: in „den absolute
% Uensartethed“ (line 2) both words are spaced; in „denne absolute
% Uensartethed“ (line 4) only „absolute“ is. The whole closing sentence „Men
% selve Ordet er undgaaet i Prof. N.'s Skrift.“ is spaced, „Men“ included --
% spacing.py flags only the run „selve Ordet undgaaet“, the short words being
% below its threshold.

Et \emph{Svar} paa Angrebet har Prof. N. altsaa ikke givet, forsaavidt er han i
god Samklang med sin Erklæring. Men hvad er der da saa gjort i Piecen paa de 52
Sider?

Der er for det Første gjort et Forsøg paa at benægte Rigtigheden af mit Referat
af Nielsens Theorie, som jeg tilmed ikke gjerne maa kalde Theorie. Naa, om
Navnet vil jeg ikke stride; kun kommer jeg lidt i Forlegenhed for at finde en
anden Betegnelse. At jeg ikke kan kalde det, der \emph{ikke} maa hedde Prof.
N.'s \emph{Theorie}, Prof. N.'s \emph{Praxis}, det har jeg til Evidents
godtgjort i min Bog (S. 214 f.). Men vi kunde jo kalde det Prof. N.'s
\emph{Mening}, thi Meningen er jo et „\emph{neutralt}“ Udtryk; eller ogsaa kunde
vi kalde det Prof. N.'s \emph{Mysterium}, hvilket af den Grund kunde være
passende, at han nødig vil ud med Sproget om det, nødig vil „spilde et Ord“ paa
at belyse det, men gjerne indhyller det i Tvetydighedens Drapperier. Dog, lad
det blive ved „Meningen“! Altsaa, Prof. N.'s Mening har jeg ikke refereret
correct, jeg har kun af Professorens Leilighedsskrifter og af
„\emph{en Andens}“\footnote{% printed *)
  \emph{\textit{Nomina}} \textit{sunt odiosa!}}
% „Nomina“ is letterspaced antiqua, „sunt odiosa!“ tight antiqua; encoded
% \emph{\textit{…}} and \textit{…} per the house convention of det-religioese.
% This is one of the antiqua phrases that carries NO guillemets.
Referater sammensat mig „et Noget“, der skal være hiin Mening, og ikke oppebiet
den „authentiske og sammenhængende“ Fremstilling. At nu denne „Andens“ Referater
ere authentiske i \emph{stræng} Forstand, vil Prof. N. vistnok ikke benægte: han
veed, at jeg
% --- p. 6 ---
\setcounter{footnote}{0}%
kan føre Beviset derfor. Men maaskee ligger Feilen i, at de ikke ere
„sammenhængende“? De ere dog givne i systematisk Form, og deres Korthed er netop
deres Fortrin; thi Meningen viser sig i dem, berøvet Phrasernes Vattering og
Tvetydighedernes Drapperier, i sin naturlige og nøgne Skikkelse, og det „Vinde
og Skjæve“ kan ikke skjules.

Det, jeg skal have misforstaaet i Professor N.'s „Mening“, skal være hans
\emph{Troesbegreb} og hans Opfattelse af \emph{Underet}. Med Hensyn hertil er
det i en mærkelig Grad undgaaet Prof. N. hvad der ikke er undgaaet mig, at
Velleiteter ikke ere tilstrækkelige. Jeg har, i Oversigten over Prof. N.'s
Mening, udtrykkelig fremhævet (S. 127) at han \emph{søger} at faae Erkjendelse
og Tanke ind i Villien, til hvilken Troen fører tilbage; men jeg har i Kritiken
af samme Mening paaviist (f. Ex. S. 196. 199), at denne Stræben ikke bliver til
Mere end en Velleitet, fordi den \emph{absolute Uensartethed} gjør hiin
Tilstedeværen af Erkjendelsen i Villien til en tom Phrase, et blot og bart
Postulat\footnote{% printed *)
  Det lykkes heller ikke bedre for Prof. N. i Piecen. Man jævnføre navnlig de
  forvirrede Bestemmelser S. 16 og 17.}.
% „Postulat“ is NOT letterspaced (spacing.py scores it 1.34; the image shows the
% sorts tight). „U n d e r e t“ in the next sentence is.
Og paa samme Maade forholder det sig med Opfattelsen af \emph{Underet}. Der
gjøres et Tilløb til, at faae en Lovmæssighed ind: den undergjørende Guds Villie
skal paa een Gang være Eet med Lovene og hævet over Lovene, og ved dette Tilløb
vilde man saa gjerne holde den slemme Conflict, Conflicten mellem Tro og Viden,
borte. Men den undergjørende Gud er for N., som jeg har viist,
„Vilkaarlighedens Gud“, „den Gud, der \emph{skaber af Intet}“\footnote{% printed **)
  \emph{Nielsen}: Om theor. og prakt. Erkj. S. 70. 74 f.}, og denne af Intet
skabende Gud ophæver Lovene ligesaa vilkaarligt som han har givet dem; thi de
ere for ham kun Love i det, der er et Intet, altsaa selv et Intet.
% There is NO closing „“ after „skabende Gud“ at the foot of that line; the
% Fraktur OCR invents one. What stands there is a speck of dirt above the d.
Det er kun et tomt Ord, naar der siges (S. 52): „det for Miraklet
Charakteristiske er slet ikke, at der gjøres et Brud paa Lovene, men det, at en
Kjendsgjerning \emph{bliver} umiddelbart \emph{henført}\footnote{% printed ***)
  % the only ***) in the book
  Man sammenligne: „\emph{det, hvori den Troende seer Miraklet!}“ (Nielsen: den
  gode Villie, S. 33).} til Guds almægtige
% --- p. 7 ---
Villie“; thi netop denne \emph{umiddelbare} Henføren til den \emph{almægtige,
skabende} Villie udtrykker en Overspringen af Lovene, et \emph{Brud} paa Lovene.
% The quotation „det for Miraklet Charakteristiske …“ opens on p. 6 and closes
% here, at the top of p. 7.

Altsaa, min Opfattelse af Prof. N.'s Troesbegreb og Underbegreb er fuldkommen
correct; og der er hverken overseet hvad Professoren \emph{vilde} eller hvad der
standsede denne Villen. Og hvad den Indsigelse angaaer, som Prof. N. gjør mod
Maaden, paa hvilken jeg har forstaaet hans Bestemmelser om Miraklets objective
Virkelighed --- dette for Prof. N. saa høist generende Punkt, der strax bringer
ham i Affect ---, saa behøver jeg kun at henvise til, hvad Prof. N. selv har
havt den Godhed at citere af min Bog. Jeg har viist, hvorledes Prof. N.'s
Forudsætninger uundgaaelig maatte drive ham til Forsøget paa, forat forebygge
Conflicten mellem Tro og Viden, at \emph{fordobble} Virkeligheden, saa at Tro og
Viden fik \emph{hver sin} Virkelighed. Men derved kom saa Prof. N. i den
ubehagelige Situation ligeoverfor de Troende at maatte borttage den objective,
ydre Virkelighed, som de fordre for Troens Gjenstand, og som de ikke kunne
opgive naar de ikke ville opgive Religionens historiske Grundlag. Prof. N. stod
saaledes mellem Scylla og Charybdis. Blev den objective Virkelighed „hvor den
virkelig udgjør et uundværligt Moment i Troesphænomenet“, fastholdt, saa blev
Conflicten med Viden uundgaaelig, og forvandledes den objective Virkelighed,
f. Ex. af Christi Fødsel, Død og Opstandelse, til en blot indre Virkelighed, saa
protesterede Troen. Thi denne Tro veed meget vel „at gjøre Forskjel paa ydre og
indre Virkelighed“, den er ikke forvirret af Phraser, og jeg har vel tydelig nok
viist, hvilken himmelviid Forskjel der ogsaa i denne Henseende er mellem den
\emph{virkelige Tro}: den Tro, der er levende i Gjerninger, og
\emph{Phantasietroen}: den Tro, der er livlig i Declamationer. Prof. N. stod og
staaer altsaa mellem to Farer, og kan kun redde sig paa samme Maade som
Blæksprutten redder sig: ved at gjøre det sort omkring sig. Jeg skal følge
Professor N.'s Exempel, og henvise hans Læsere til Christiania\-%
% --- p. 8 ---
% „Christianiaforelæsningerne“ is broken by the compositor at the foot of p. 7
% and rejoined here with \-% so that the transcription reflows.
forelæsningerne IX og X, og saa kun bede dem, med lidt Eftertanke at gjennemlæse
Svaret til Prof. Caspari i Efterskriften; de ville saa forstaae, hvorledes det
Middel anvendes: „at tale sort“.
% The roman numerals „IX og X“ are ANTIQUA in the print, as all figures and
% roman numerals in this book are; set plain here, not \textit{}.

Prof. N.'s Indsigelse mod min Opfattelse af hans Mening kunne vi vel saa betragte
som affærdiget. Men naar den er fjernet, saa har Professoren en ny Vending paa
rede Haand, og det en saadan, over hvilken han synes at være saare glad: Ved min
Kritik af „Meningen“ skal jeg netop kun have opnaaet hvad den snilde Strateg, mod
hvem jeg kæmpede, selv ønskede: at stadfæste, hvad Prof. N. atter og atter
indskærper, at Troeslæren umulig kan bestaae for Videnskabens Domstol, at den
derfor hverken kan eller vil underkaste sig Videnskabens Afgjørelse. I min
Ubesindighed er jeg altsaa falden i det snildt lagte Baghold, og min snilde
Modstander synes saare glad. Det turde dog maaskee ikke være saa ganske
paalideligt med denne Glæde; thi Prof. N. seer sig nødsaget til at gjøre nogle
smaa Behændighedskunster forat overbevise Andre om, at den er begrundet. Det,
som jeg har rettet min Kritik imod, er Meningen om \emph{Troens og Videns
Forhold}, et Forhold, som, efter Prof. N.'s egen Bestemmelse, Viden erkjender,
idet den, medens den begriber sit eget Væsen, afspeiler Modsætningen. Min
\emph{Videnskritik} er altsaa rettet mod \emph{Vidensbestemmelser}, min Taktik
er altsaa, efter Forudsætningerne, aldeles correct. Men forat den ikke skal være
det, lader Prof. N. først det, jeg ikke maatte kalde Theorie, blive til „den af
mig (ɔ: Prof. N.) opstillede \emph{Troestheorie}“ --- altsaa dog
\emph{Theorie}! --- og saa, ved en lille hurtig Manøvre, \emph{Troestheorien}
blive til \emph{Troeslæren}.
% ɔ: -- the id-est mark (reversed c + colon), U+0254 + colon, verified at
% 400 dpi. RECON.md §8 records it only for p. 13; this is a second occurrence.
% The Fraktur OCR reads it „o:“ and the ABBYY layer „໑:“.
Og saa er jeg naturligviis overbeviist om at have ladet mig fange. Men selv om
vi nu lode denne Substitution staae, saa har jeg jo netop i min Bog godtgjort,
hvorledes Dommen over alt det, der bliver et Udtryk for det menneskelige Væsen,
i sidste Instants tilkommer Viden ifølge dennes Forhold til Væsenet; jeg har
ikke postuleret det, som N. paastaaer (Piecen S. 8), men ført „Grund\-%
% --- p. 9 ---
% „Grundbeviset“ is broken at the foot of p. 8 and rejoined here.
\setcounter{footnote}{0}%
beviset“ baade i \emph{directe} og \emph{indirecte} Form og godtgjort
Vidensbevisets Gyldighed\footnote{% printed *)
  Man see i min Bog de Udviklinger, der sammenfattes S. 210 f. jfr. S. 169.}.

Prof. N.'s Glæde over den snilde Vending var altsaa kortvarig, han maa nu dog
til at iværksætte Tilbagetoget, og det skal saa kun dækkes ved en lille offensiv
Bevægelse. Ved den bliver det først ganske interessant, at Prof. N. kommer til
at begaae en lignende taktisk Feil, som den, han ovenfor vilde paadutte mig. Om
nemlig ogsaa det „offensive Stød“ var faldet nok saa heldigt ud, om han havde
tilintetgjort Alt, hvad jeg gjennem min Kritik havde givet af positive Bidrag
til en psychologisk og metaphysisk Theorie, saa vilde Prof. N. endnu være lige
nær. Han vilde have staaet ligeoverfor Umuligheden af at fastholde den
„\emph{absolute Uensartethed}“, denne Bestemmelse, der bærer hele Prof. N.'s
Theorie --- nu kunne vi jo atter kalde den \emph{Theorie}! --- og med hvilken
den staaer og falder. Og naar han saa opgav den --- hvad han maaskee nu ikke var
saa utilbøielig til, naar det kunde skee paa en lempelig Maade, --- og fik sig
en ny Theorie sat sammen, der skulde føre til samme Maal som den gamle, saa
vilde han staae ligeoverfor det uløselige Problem, hvorledes der, naar Videns
normative Gyldighed for Personlighedens Helhed benægtes, kan findes en Norm, der
paa det Ethiskes og Religiøses Gebet gjør Skilsmisse mellem det Sande og det
Usande, mellem Personlighedens sande og væsentlige Fordring og Individets usande
Stræben\footnote{% printed **)
  Jfr. mit Skrift f. Ex. S. 169. 203.}. Jeg har i mit Skrift paaviist
Umuligheden deraf: Prof. N. vil ligesaalidt som nogen Anden kunne godtgjøre
Muligheden.

Hvorledes end altsaa Prof. N.'s Angreb var faldet ud, vilde hans Theorie lige
lidet være reddet; mit Skrifts negative Resultat vilde staae urokket. Men Prof.
N.'s Angreb er saa lidet farligt, at det med ringe Møie kan kastes tilbage, ja
det er endog i det Væsentlige allerede affærdiget i det Skrift, mod hvilket det
er rettet.

% --- p. 10 ---
% p. 10 carries only SEVEN lines of body text; the rest of the leaf (some 33 of
% its 40 lines) is the single note *), which then runs over onto p. 11.
\setcounter{footnote}{0}
Af selve Angrebet bliver det naturligviis kun Hovedpunkterne, jeg kommer til at
berøre; at gaae ind paa de enkelte smaa Vendinger, smaa Vittigheder og smaa
Slutninger med store Paralogismer, vil være ganske overflødigt ligeoverfor
enhver Læser, der ikke enten af Naturen er meget uklar eller er bleven det ved
uforsigtig Nydelse af den Art Logik, der gives paa de „60 Ark“\footnote{% printed *)
  % THIS NOTE RUNS OVER THE PAGE BREAK IN THE PRINT: it fills the note area of
  % p. 10 and its last two lines („To ville, „paa Grund af“ … blive til Een.“)
  % stand unmarked in the note area of p. 11, which carries no note of its own.
  % It is set here IN FULL at its call, per RECON.md §4; see the % comment at
  % the p. 11 page marker below.
  Saaledes mener jeg, trøstig at kunne overlade til enhver Læser, der ikke lider
  af ovennævnte naturlige eller erhvervede Svaghed, selv at komme tilrette med
  hvad Prof. N. udvikler S. 8 f. om \emph{Existents og Theorie} og S. 20 f. om
  \emph{Villiens Forhold til Selvet og dets Selvbestemmen}. Man behøver blot at
  passe paa, hvad Prof. N. forsigtig lister sig forbi, og hvad han, ikke meget
  behændigt, lister ind, for at være orienteret. Ligesaa rolig troer jeg at
  kunne overlade ham Kritiken af de forbausende Consequentser, Prof. N.
  (S. 24 f.) faaer ud af hvad jeg har sagt om Enheden af en \emph{bestemt},
  særegen Form af den theoretiske Tænkning, og den Tænkningens Form, der vilde
  være den, efter Forudsætningen med al Theorie \emph{absolut uensartede},
  % „ab=/solut uensartede“ is broken across a line in the print and the
  % letterspacing carries across the break; rejoined here.
  „praktiske Tænknings“. Det vilde vel heller ikke være ham uoverkommeligt, selv
  at korrigere hvad Prof. N. slutter af, at en Personlighed aldrig fuldstændig
  kan blive gjennemsigtig for en \emph{anden} Personlighed, eller at indsee,
  hvilken Forskjel jeg gjør mellem Prof. N.'s „absolut uensartede“ Tro og Troen
  i Religionernes Betydning. Endelig turde det vel ikke heller falde ham
  vanskeligt, ved det „Exempel \textit{instar omnium}“ paa mine „tomme Ordspil“,
  over hvilket Prof. N. er saa henrykt (S. 50), at gjøre den Opdagelse, at
  Eftertrykket falder paa, at Mysteriet, der, efter Prof. N.'s Mening skulde
  \emph{forhindre} Troens og Videns gjensidige Indskrænkning, netop kommer til
  at \emph{udtrykke} Indskrænkningen, og at det „Fine“ ved, at „Mysteriet“
  \emph{ikke} bliver et Mysteriøst, er det, at en Hemmelighed ikke længer er en
  Hemmelighed, naar man veed \emph{hvad dens Indhold er}. Men netop det vides,
  naar Tro og Viden ere hinandens Mysterium, og ere tilstede i een og samme
  Bevidsthed. Men her har rigtignok den „absolute Uensartethed“ --- denne Prof.
  N.'s store og fatale Opdagelse --- spillet hans „gode Villie“ det Puds, at den
  Bevidsthed, der skal rumme baade Tro og Viden, virkelig er „spaltet“. Man læse
  i „den gode Villie“ (S. 33): „idet den Troende og den Vidende“ (der vel at
  mærke er \emph{samme} Individ) „\emph{begge} paa een Gang rette deres Blik mod
  Mysteriet, udlignes Striden og Extremerne enes“! --- Altsaa, ligesom i
  Elverdandsen den Ene blev til Tre, bliver her den Ene til To, men de To ville,
  „paa Grund af“ den absolute Uensartethed, desværre ikke atter blive til
  Een.}.
% --- p. 11 ---
% The note area of p. 11 holds ONLY the unmarked last two lines of p. 10's *)
% („To ville, „paa Grund af“ den absolute Uensartethed, desværre ikke atter
% blive til Een.“). They are set above, at the call on p. 10, and are NOT
% repeated here; p. 11 has no note of its own and no \setcounter is needed.
Min første Synd skal være den, at jeg confunderer det Almindeliges videnskabelige
% RETTELSE. The author's own erratum leaf (PDF 37) reads „Side 11, Lin. 2:
% videnskabelige, læs: subjective.“ and it is right: the sense wants
% „subjective Gyldighed med den objective“. Per the repo's diplomatic stance the
% PRINTED reading stands here and the Rettelse is set at the end of the edition;
% it is NOT applied.
Gyldighed med den objective, endog i den Grad, at jeg ikke seer mig istand til
at gjøre Forskjel mellem Videnskabsgrunde og Samvittighedsgrunde, mellem en
Erkjendelse af Pligtloven og en Erkjendelse af Tyngdeloven. Og denne min Synd er
den svareste, hvorfor ogsaa Prof. N. ender sin Piece med at kræve mig til
Regnskab for den. Men her er jeg nu saa heldig, allerede at have aflagt
Regnskabet, hvilket kan eftersees i min Bog, navnlig Side 141---148. Kun med
Hensyn til min Mangel paa Evne til at gjøre en Forskjel imellem en Erkjendelse
af Pligtloven og af Tyngdeloven maa jeg gjøre opmærksom paa, at denne
Sammenstilling er Prof. N.'s, ikke min. Jeg har sammenstillet den theoretiske
Erkjendelse af Menneskets Væsen og Erkjendelsen af den Pligt, der er bestemt ved
Væsenet, som bliver sig selv aabenbart i hiin theoretiske Erkjendelse, og paa
Grund af Pligtens Forhold til Væsenet har jeg fastholdt Enheden af theoretisk
Erkjendelse og \emph{Pligterkjendelse}, den \emph{theoretiske} Erkjendelses
Betydning som nødvendig Betingelse for Handlingens Sandhed. Naar Prof. N. i sin
„Sagudvikling“ ikke havde benyttet den Fremgangsmaade, som Dr. Heegaard saa
skarpt og træffende betegner i sit lille Skrift, S. 29 f., men istedenfor den
objective Erkjendelse af Tyngdeloven, der „slet ikke angaaer Personlighedens,
men kun Legemlighedens Verden“, havde sat en objectiv Erkjendelse af hvad der
falder indenfor det Psychologiske: af selve det menneskelige Selv, den
menneskelige Personlighed og dens Grundfunctioner, saa vilde den Erkjendende
ikke have været „saa udtrykkelig opfordret til at see bort fra al
Personlighed“, og saa vilde „Sagudviklingen“ maaskee have ført til et andet
Resultat.
% „Dr.“ is antiqua here as on p. 3; left plain, see the comment at p. 3.

Min anden Synd skal være den, at jeg af en Feiltagelse skal sætte
Enkeltvilliens Almeenbegreb i den for sin Selvstændighed kæmpende Villies Sted.
Maa jeg gjøre Prof.
% --- p. 12 ---
N. et Spørgsmaal? Naar der skal bestemmes, hvad den enkelte Villie \emph{som
enkelt} gjør idet den bestemmer sig til i Handling at virkeliggjøre det
Erkjendte, hvorledes skal der saa kunne \emph{tales} om denne Villiesact, naar
der ikke tales i \emph{Begreber}, der som Udtryk for Enkeltvilliens Væsen
udtrykker Enkeltvilliens Lov?
% PRINTER'S DEFECT, p. 12, line 4 of the page: a solid round blot of ink stands
% in the word space between „Væsen“ and „ud-/trykker“ (verified at 400 dpi and
% at 6x). It is foul matter or a bearer taking ink, not a hyphen -- the
% line-break hyphen after „ud“ is separately visible -- and not a sort, so no
% sort-level claim is made from it (JBIG2; TRANSCRIPTION-PLAYBOOK §2). The two
% words are transcribed with an ordinary space. The Fraktur OCR reads the blot
% as „*“, the ABBYY layer ignores it.
Dette er ikke at identificere Enkeltvilliens Almeenbegreb med det enkelte
Individs Villie, men at sætte denne som det, i hvilket hiint Almeenbegreb faaer
Virkelighed, og som derfor har sin gyldige Lov i det. Naar derimod Prof. N. fra
sit Standpunkt taler om Villiesbegrebet, saa er der den Mærkelighed derved, at
dette Villiesbegreb \emph{ikke} bliver et Begreb om Villien, altsaa kaldes
Villiesbegreb paa samme Maade som \textit{lucus a} \emph{\textit{non}}
\textit{lucendo}.
% Antiqua, with ONLY „n o n“ letterspaced: the print gives „lucus a n o n
% lucendo.“ Encoded \textit{…} for the tight antiqua and \emph{\textit{…}} for
% the spaced word, per the house convention. No guillemets here.
Villien i Prof. N.'s Opfattelse bliver nemlig, hvormeget han end strider imod,
det, der ved sin absolute Uensartethed med Viden, ikke kan finde et Udtryk i
Begrebet. Saa maatte man da ialfald lade Fælledsbestemmelsen: Begreb omfatte to
absolut uensartede Underarter, saaledes som Underarterne af Prof. N.'s
Skabelsesbegreb og Mirakelbegreb (jfr. min Bog S. 206 f.); men et saadant
Forhold er en Umulighed, og det bliver kun et tomt Spil med Ord, at betegne det
absolut Uensartede ved samme Bestemmelse. Vilde Prof. N. kalde sit
Villiesbegreb et „Afspeilingsbegreb“, saa behøver jeg kun at henvise ham til min
Belysning af denne „Kategorie“: Afspeiling (i mit Skrift S. 148 ff.).

Vi komme saa til min tredie og sidste Synd: min sørgelige Forvirrethed i
Bestemmelsen af Selvhedsenergien og af Videns og Villiens Forhold til Væsenets
Grundbegreb. Her er det imidlertid hændt Prof. N., at han paa en meget
besynderlig Maade har misforstaaet nogle Grundbestemmelser, jeg har givet,
Bestemmelser, som jeg aldrig har havt Vanskelighed ved at faae forstaaede af de
yngste Studerende, og som Hr. Professoren vist ogsaa let vilde have opfattet
rigtigt, hvis han ikke havde fordybet sig saaledes i sin egen Terminologie, at
han kun ret forstaaer hvad der udtrykkes i den.

Jeg har, i de psychologiske Udviklinger i mit Skrift, fremhævet
Væsensuniversaliteten som det, der i Modsætning
% --- p. 13 ---
til Dyrevæsenets Particularitet: en Particularitet, der røber sig i Instinctlivet,
% PRINTER'S DEFECT, p. 13, line 1: a thin vertical stroke, dotted at the top,
% stands in the word space immediately before the second „Particularitet“ and
% touches the P (read at 600 dpi). Foul matter or spacing material taking ink --
% not a sort and not punctuation, and no sort-level claim is made from this
% JBIG2 scan (TRANSCRIPTION-PLAYBOOK §2). Transcribed as an ordinary word space.
% The Fraktur OCR reads it „¡“; the ABBYY layer ignores it.
charakteriserer Menneskevæsenet. I denne Væsensbestemmelse, der indeslutter den
indre Uendeligheds Bestemmelse, har jeg søgt Selvbevidsthedens Mulighed,
Muligheden altsaa til det Selvets Forhold til sig selv, til den Selvets
Forsigværen, der, i Modsætning til det dyriske Selvs Forsigværen i
Selvfornemmelsen, i hvilken det Almene vel er tilstede, men altid som bundet til
den \emph{enkelte} Fornemmelse, naturlig kan siges at have den frigjorte ɔ: den
% ɔ: -- the id-est mark (reversed c + colon), U+0254 + colon, verified on the
% image at 600 dpi. RECON.md §8 records this as the one confirmed instance in
% the book; no second one was found in pp. 13--22. tesseract gives „o:“ and
% „9:“; the ABBYY layer is the only witness that reads it at all.
fra denne Bundethed frigjorte, Almindeligheds Form; og ifølge Erkjendelsens
Forhold til Væsensuniversaliteten (S. 137) har jeg tillagt den det ideelle
Principat i Menneskeaanden.

Prof. N. har nu paa en besynderlig Maade misforstaaet Betydningen af
Bestemmelsen: Væsensuniversalitet, og af Betegnelsen: den frigjorte
Almindelighed, og i Kraft af disse Misforstaaelser drager han en Mængde
Slutninger af en ganske eiendommelig Art (S. 26 ff.), gjennem hvilke der
paademonstreres mig Absurditeter. Men saasnart Misforstaaelserne fjernes, saa
tage de Absurditeterne med sig.

Saa har jeg, i mit Skrift S. 135, bestemt Erkjendelsens og Villiens Forhold til
„det i Selvbevidstheden som Enhed aabenbarede Væsens Energie“, og for de
Bestemmelser, jeg der har givet, maa jeg høre meget Ondt af min ærede Kritiker.
Først skal jeg have laant mine „Slagord“: Selvhed og Selvhedsenergie fra
Prof. N.'s Philosophie. Det er mig lidt overraskende. Ikke blot har jeg brugt
disse Kategorier i psychologiske Forelæsninger, der ere holdte før noget af de
Skrifter af Prof. N. var udkommet, i hvilket han har \emph{berørt} psychologiske
Spørgsmaal, --- noget sammenhængende Psychologisk har Prof. N. aldrig givet. Men
jeg stod ogsaa i den naive Formening, at Selvhed --- det danske Ord for
Subjectivitet --- var en courant Kategorie i den efterkantiske Philosophie: en
% „courant“ is Fraktur here, not antiqua; so is „à la“ on p. 18. Only Latin and
% Greek are set in antiqua in this book (RECON.md §5).
Kategorie, som hverken Prof. N. eller jeg har opfundet eller behøvede at opfinde.
Og hvad Selvhedsenergie angaaer, et Udtryk, jeg forresten ikke mindes at være
stødt paa hos Nielsen, der ialfald, ifølge den absolute Uensartethed, ikke vilde
kunne gjøre Stort ud
% --- p. 14 ---
deraf, saa er det Selvhedens Bestemmelse, forbunden med den gamle Bestemmelse af
ἐνέργεια, der skyldes Psychologiens Grundlægger Aristoteles. Fra Aristoteles's
% Greek, antiqua italic. The breathing and the accent were read off the image at
% 600 dpi and at 8x: a smooth breathing over the first epsilon and an acute over
% the second -- ἐνέργεια. The Fraktur model reads no Greek at all; the ABBYY
% layer reads the word but its diacritics are not to be trusted (RECON.md §8).
Skrifter har jeg hentet denne Bestemmelse, ikke fra Prof. Nielsens, og skjøndt
der i disse undertiden røber sig en paafaldende ubevidst Tilbøielighed hos deres
Forfatter til at annectere hvad der tilhører ældre Philosopher, saa kan jeg dog
% PRINTER'S DEFECT, p. 14: a short horizontal stroke stands high in the word
% space between „saa“ and „kan“, about at ascender height (600 dpi). Ink on
% spacing material or foul matter, not a sort; transcribed as an ordinary word
% space. tesseract reads it as a closing quote, „saa” kan“; ABBYY ignores it.
neppe forestille mig, at Prof. N. skulde gaae i den Selvskuffelse, at han,
maaskee i sin Ungdom, paa et af sine mange tidligere Udviklingsstadier, skulde
have skrevet den Psychologie, der saa i længere Tid har gaaet under en vis
Aristoteles's Navn. Om Prof. N. kan jeg, som sagt, ikke forestille mig det, men
% PRINTER'S DEFECT, p. 14: a round ink dot at mid-height stands in the word
% space before „Navn“ (600 dpi). Same class as the two marks above and as the
% blot logged on p. 12; transcribed as an ordinary word space. tesseract reads
% it „‘Navn“, ABBYY „" Navn“.
for hans Tilhængere --- eller Tilhænger; thi nu bør man maaskee tale i
Enkelttallet, --- tør jeg ikke svare; han har jo allerede opdaget kjendelige Spor
af sin Mester i den nordiske Mythologie.

Hine „Slagord“, som jeg altsaa \emph{ikke} har laant fra Prof. N., skal jeg saa
have handlet ilde med idet jeg nærmere bestemte dem. Jeg skal have faaet en
besynderlig Trehed af Selvhedsenergier ud, hvis Forvirring i Prof. N.'s
Beskrivelse (S. 38 ff.) næsten tager sig ligesaa forfærdelig ud som Forvirringen
blandt hans Grundideer. Men et Par korte Bemærkninger turde maaskee vise, at hele
Prof. N.'s Consequentsdrageri er begrundet paa Misforstaaelser. Hvad jeg har sagt
er Følgende: Det menneskelige Væsens Grundbestemmelse er Selvhedsenergie; i
Selvbevidsthedens Enhed aabenbarer Væsenets Enhed sig, og for dette udelte
Væsens Energie er Erkjendelse og Villie primitive Udtryk, primitive
Realisationsformer. Selvhedsenergiens almene Grundbestemmelse faaer altsaa sine
særlige Udtryk i Erkjendelsesenergiens og Villiesenergiens Former; i begge disse
primitive Former udtrykker Væsenet sig saaledes, at dets Enhed bevares --- ikke,
som Prof. N. vil, \emph{adsplittes} i absolut Uensartethed --- i begge Former er
derfor Indholdet \emph{væsentligt} det samme; thi Indholdet er
Væsensbestemmelse, og Formernes abstracte Modsætning ophæves idet de
Bestemmelser, der udtrykke den, vise hen paa hinanden og udvikle sig af hinanden
(jfr. mit
% --- p. 15 ---
\setcounter{footnote}{0}
Skrift S. 134 f.). Grundformerne maae, netop som Grundformer, reflectere
hinandens Formbestemtheder.

Naar man sammenholder disse Sætninger med hvad Prof. N. med synlig
Tilfredsstillelse og utallige Gjentagelser har skrevet i sin Piece S. 36 ff., saa
vil man see, at Hr. Professoren med lidt mindre Hidsighed og lidt mere
Eftertanke kunde have sparet sig selv adskillige Ikke-Kløgtigheder.

Skulde Læseren med Hensyn til et enkelt Punkt, f. Ex. med Hensyn til
Erkjendelsens og Villiens Forhold paa de ufuldkomne Udviklingstrin og paa deres
Culminationstrin, eller til Erkjendelsesindholdets \emph{Omsætning} i
Villiesindhold og omvendt, ønske nærmere Oplysning, saa henviser jeg ham til mit
Skrift (S. 136 ff.), og navnlig til det, Prof. N. har udeladt i sine
Citater\footnote{% printed *)
  Prof. N. citerer idethele noget skjødesløst. Saaledes er det Citat, der findes
  i hans Piece S. 22. Lin. 6---11, givet i en aldeles meningsløs Form, og
  Forvanskningen er tilmed gjentagen med Citationstegn længere nede paa samme
  Side.}.

Ved Siden af de paafaldende Misforstaaelser, Prof. N. har gjort sig skyldig i med
Hensyn til de anførte Punkter, faae de Beskyldninger for Misforstaaelse, han
retter imod mig, en eiendommelig Belysning. Jeg har allerede paaviist, hvorledes
det hang sammen med min formentlige Misforstaaelse af Prof. N.'s Troesbegreb og
Underbegreb; og med et Par Ord vil jeg kunne paavise, hvorledes andre formentlige
Misforstaaelser opløse sig i fuldkommen rigtig Forstaaelse. Jeg skal, efter
Piecen S. 32, i min Opfattelse af Villien have gjort mig skyldig i en saa
mærkværdig Misforstaaelse, at Ingen der „fornuftigviis maatte antages at kunne
læse“ skulde være istand dertil. Misforstaaelsen skal bestaae deri, at jeg har
ført Prof. N.'s Villiesbegreb (\textit{sit venia verbo!}) tilbage til den
% ANTIQUA, verified at 600 dpi: „sit venia verbo!“ is upright antiqua standing
% against the Fraktur, like the other Latin in this book, and carries no
% guillemets. It is NOT in RECON.md §5's list of antiqua phrases -- that list is
% therefore not exhaustive; this is a new find of the pp. 13--22 batch.
„umiddelbare Naturbestemtheds Virksomhedsyttring i Driften, den sandselig
egoistiske Attraa“, og det sætter Prof. N. sig med Hænder og Fødder imod. Men,
som jeg ovenfor har viist, Velleiteter hjælpe ikke, og med den absolute
Uensartethed mellem Viden og Villie til For\-%
% --- p. 16 ---
% „Forudsætning“ is broken by the compositor at the foot of p. 15 and rejoined
% here with \-% so that the transcription reflows.
udsætning, kan Villien, i sin Sondring fra Videns Almene, virkelig ikke blive til
Andet end hvad jeg har ladet den blive til. Det nytter sandelig ikke, forat redde
en Erkjendelse og en Tænkning, og derigjennem det Almenes Bestemmelse, for denne
Villie, at statuere to \emph{absolut uensartede} Sorter Erkjendelse og Tænkning;
thi vi have allerede seet, at det kun er at gjøre Bestemmelsen Erkjendelse og
Tænkning til et tomt Ord. Altsaa her er ingen Misforstaaelse paa min Side; men
Prof. Nielsen har kun ikke forstaaet, til hvilke Consequentser den absolute
Uensartethed --- denne absolute Uensartethed, som han nu tier saa stille om ---
uundgaaelig maa føre.

Et andet Sted (S. 12) meddeler Prof. N. sine Læsere, at jeg har gjort mig skyldig
i „Mere end en Fordreielse“ af hans Ord, i „en gjennemgaaende Misforstaaelse af
den hele Sag, hvorom der handles“. Denne Sag er Forholdet mellem Subjectiviteten
og Objectiviteten. Professoren tager meget stærkt paa Veie og bruger adskillige
Sider til at belære og revse mig. Og hvad er saa det Hele? En Misforstaaelse,
ikke fra min, men fra Prof. N.'s Side. Jeg havde som en Slags Indledning til
Kritiken af Prof. N.'s Theorie belyst Forvirringen i hans Opfattelse af nogle af
de Bestemmelser, ved hvilke han udtrykte Grundmodsætningen: Bestemmelserne
Existents og Theorie, Subjectivitet og Objectivitet. Jeg vilde deels gjøre det
tydeligt, hvorfor jeg i Kritiken holdt mig til et andet Udtryk for Modsætningen,
deels vilde jeg fra en ny Side belyse Forholdet mellem Kierkegaard og R. Nielsen
ved at vise, hvilken Skjæbne den fine og skarpe Tænkers klare Kategorier fristede
ved at komme i Prof. N.'s Hænder. Jeg paaviste saa Forvirringen og Tvetydigheden
i Prof. N.'s Bestemmelser, og den sophistiske Brug, han gjorde af disse. Med
Hensyn til Modsætningen mellem Subjectivitet og Objectivitet, saaledes som
Prof. N. bestemte den, gjorde jeg blandt Andet opmærksom poa to Ting. Det Ene
% PRINTER'S DEFECT, p. 16: the print reads „poa“ for „paa“. Read at 600 dpi and
% at 8x against two controls on the same page: the second letter is a closed
% round o, identical with the o of „to“ two words later, and quite unlike the
% flat-topped Fraktur a of „paa“ four lines below. Transcribed as printed, per
% the diplomatic stance; both OCR witnesses read „poa“ as well. (The scan is
% symbol-matching JBIG2, so this is recorded as a reading of the image, not as a
% claim about what the compositor's case held -- TRANSCRIPTION-PLAYBOOK §2.)
var, at det, idet den abstracte Modsætning definitivt fastholdes, vilkaarligt
bliver ignoreret, at paa de forskjellige Subjectivitets-Stadier Subjectets
Bestemthed ---
% --- p. 17 ---
\setcounter{footnote}{0}
hvad der især tydelig fremtræder paa det Stadium, der betegnes som den uendelige
Subjectivitets --- netop er i Kraft af hvad der er bestemt som Tilværelsens
objective Princip, og at den udtrykker dette. Det Andet var, at det bliver
overseet, at den subjective Sandhed netop er den tilegnede og i Tilegnelsen med
sig væsentlig lige objective Sandhed. Ved ikke at forstaae, at den sidste Passus
generelt udtrykker Forholdet, og ved, paa en høist vilkaarlig Maade, at referere
den exclusivt til et enkelt af de nævnte Stadier, faaer Prof. N. ud, at jeg
aldeles skal have misforstaaet hans Ord om dette, og at jeg skal have viklet mig
ind i gyselige Modsigelser. Saasnart mine Ord forstaaes rigtigt, seer man, at de
ingen Misforstaaelse af Prof. N. forudsætte, og at de ere fuldkommen correcte.
Det viser sig især tydeligt paa Stoicismens og Hedonismens Standpunkt, at
Subjectets Bestemthed udtrykker Tilværelsens objective Princips Bestemthed. Dette
behøver ingen Forklaring for den, der blot kjender lidt til Philosophiens
Historie, og ikke heller behøver det en Forklaring, at Tilværelsens Princip
betegnes som Tilværelsens objective Princip. Og naar saa den generelle Sætning om
Forholdet mellem den subjective og den objective Sandhed anvendes paa disse to
Systemer, saa bliver altsaa Consequentsen, at den relative Sandhed, som det
enkelte ensidige, men i sin Ensidighed et Gyldigt indesluttende System eier, er
Maalet for den subjective Sandhed hos den, der practisk tilegner sig Systemets
objective Sandhed\footnote{% printed *)
  Paa Prof. N.'s snurrige Argumentation fra det, at en Stoiker kan „sætte sig
  ligesaa sikkert ind i“ den hedoniske Theorie som en Hedoniker og dog ikke
  blive Hedoniker, skal jeg kun henlede Læserens Opmærksomhed: den behøver ingen
  Commentar!}.

Saaledes hænger det altsaa sammen med mine Misforstaaelser, og saaledes er det
Angreb beskaffent, ved hvilket Prof. N., som ved en Arrieregarde-Fægtning, søger
at dække sit factiske Tilbagetog. Men Prof. N. gjør ikke gjerne dette Tilbagetog;
han er bleven vred, saare vred, hidsig, saare hidsig. Og denne Vrede og denne
Hidsighed maa han give Luft paa
% The signature numeral „2“ stands at the foot of p. 17 and is not transcribed.
% --- p. 18 ---
\setcounter{footnote}{0}
mange Maader. Snart maa jeg høre Ondt for min „kjendelige Selvfølelse“, snart
fordi jeg „paa et luftigt Underlag af væsenstomme Abstractioner, skjæve
Deductioner, udtværede\footnote{% printed *)
  Haarde Beskyldning i den Mands Mund, der har skrevet anden Deel af
  Grundideernes Logik!} Vendinger og stakaandede Bemærkninger“ har „styltet mig
op til en philosophisk Myndighed, en Videns Ypperstepræst“. Snart bliver min
Kritik beæret med Titelen af „tom Ordkritik“ og der tillægges mig Bemærkninger
à la Per Degn. Saa meddeles der mig, at jeg mangler Agtelse for Videnskaben og
Sands for den, og saa culminerer endelig Prof. N.'s velvillige Stemning, idet han
paatager sig at søge Motiver for min Optræden. Om det sidste Punkt skal jeg siden
tale et alvorligt Ord med Hr. Professoren; Et og Andet af det Øvrige skal jeg her
i Korthed belyse; der vil med det Samme falde nogle, maaskee lidt uvelkomne,
Streiflys paa min Modstander.

Prof. N. vil reducere min Kritik til en Ordkritik, der mangler „al
Sagudvikling“. Denne Vending har Prof. N. laant fra hiin ene Tilhænger, der dog
allerede har brugt den saa ofte, først mod Cand. Brandes, siden mod Dr. Heegaard,
at den nu turde være lidt forslidt. Det „Exempel \textit{instar omnium}“ paa
% ANTIQUA inside the Fraktur quotation marks, verified at 600 dpi; the „…“ are
% Fraktur, the Latin is antiqua and carries no guillemets. Same phrase as p. 10.
„tomt Ordspil“, som Prof. N. anfører imod mig, har jeg allerede ovenfor belyst
(S. 10), og det viste sig dog at være adskilligt Mere end et Ordspil. Omvendt har
jeg belyst Prof. N.'s „Sagudvikling“ af Forholdet mellem Erkjendelsen af
Tyngdeloven og af Pligtloven, og det viste sig, at den just ikke saa ganske
udviklede Sagen. Det vilde være at slippe altfor let for Prof. N., naar han,
hvergang hans Begrebsudviklinger bleve paaviste som falske og forvirrede, blot
kunde raabe „Ordkritik“ og saa „almindeligt i al Almindelighed“ pukke paa
Sagudvikling. Tankerne haves kun i Ordene, og det maatte være meget vage og løse
Tanker, der ikke kunde taale at fastholdes i den Ordets Form, i hvilken de
fremtræde. Og Sagudviklingen er just ikke Prof. N.'s stærke Side. Man lade sig
blot ikke blænde af, at Prof. N.
% --- p. 19 ---
har den Skik, i sine Skrifter at optage store Stofmasser fra de heterogene
Videnskaber, med hvilke han sysler som Dilettant, og man lade sig endnu mindre
blænde af, at han taler pompøst om „Stofforbrænding i Dialektikens Ildluft“
o. s. v. „Stofforbrændingen“ er saare liden, og ved lidt nærmere Eftersyn vil det
vise sig, at der i de allerfleste Tilfælde ikke bliver mere organisk Sammenhæng
mellem Prof. N.'s Begrebsudviklinger og det Stof, der tages med, end i et
zoologisk Museum mellem Dyrehuderne og de Blaar, hvormed de ere udstoppede. Men
der gives ganske vist dem, hvem der er sat saa mange Blaar i Øinene, at de ikke
kunne see det!

Saa gjør Prof. N. mig til en anden Per Degn! Naa ja! Rasmus Berg kunde jo gjøre
Per Degn til en Hane, hvorfor skulde saa Rasmus Nielsen ikke kunne gjøre mig til
Per Degn? Jeg skulde saamænd heller ikke have spildt et Ord derpaa, hvis ikke
Prof. N., for at faae Leilighed til at anbringe sin Vittighed, paa en noget
aparte Maade havde omgaaedes med Citater. Naar Prof. N., istedenfor at have
begyndt Citatet af sin egen Bog med Ordene: „Staaer det fast, at Troen o. s. v.“,
havde taget de Meningsløsheder med, der gaae forud, og naar han, istedenfor at
begynde Citatet af min Bog med: „De livlige Analogier o. s. v.“, havde taget det
Forudgaaende med (S. 187), i hvilket hine Meningsløsheder ere paaviste, saa vilde
han have orienteret Læseren betydeligt bedre. Og naar han havde citeret den af
mig tilføiede Anmærkning fuldstændigt, saa vilde Læseren strax have kunnet see,
fra hvem den Per-Degnske Oplysning kommer.
% The print sets „Per = Degnske“ with the Fraktur hyphen and a word space on
% either side of it; transcribed „Per-Degnske“, the hyphen kept as printed
% (RECON.md §8).

Jeg havde i Anmærkningen, i Anledning af nogle \textit{«lazzi»} af Prof. N.,
% ANTIQUA with ANTIQUA GUILLEMETS, verified at 600 dpi. The guillemets follow the
% fount, not the sense: most of the antiqua Latin in this book carries none
% (RECON.md §3). Both marks are set inside \textit{} because both are antiqua.
sagt: „Nielsens Exempler med de „mathematiske Venner“ og den „venskabelige
Mathematik“ oplyse kun det lidet Interessante, at man ikke uden videre, fordi et
Adjectiv kan prædiceres om et Substantiv, kan prædicere et til Substantivet
svarende Adjectiv om et til Adjectivet svarende Substantiv; at man f. Ex. ikke,
fordi man kan tale om en flau Vittighed, kan tale om en vittig Flauhed. Det,
% Single-level nesting, as printed: the outer „…“ opens here and closes on p. 20
% after „gjennem dem.“, and the inner pairs round „mathematiske Venner“ and
% „venskabelige Mathematik“ are ordinary „…“, not the doubled „„…““ of p. 20.
% The signature numeral „2*“ stands at the foot of p. 19 and is not transcribed.
% --- p. 20 ---
at der kan gives ethiske og uethiske Astronomer, oplyser kun, at forskjellige
Aandsvirksomheder kunne have samme Subject uden at kunne gjøres til Prædicat for
hinanden (saaledes som f. Ex., indenfor Videns Sphære, Astronomie og
Lægevidenskab, idet der vel gives lægekyndige Astronomer, men ingen
lægevidenskabelig Astronomie); men det beviser Intet med Hensyn til en absolut
Uensartethed, eller til, at disse Aandsyttringer kunde prædiceres om det samme
Subject, uden at de væsentlig gik tilbage til samme Princip, eller uden at det
samme udelte Væsen gjorde sig gjældende gjennem dem.“ Hos Prof. N. læser man nu
(S. 50. Anm.) Følgende: „Istedenfor at oplyse Noget om Sagen, har Hr. Prof.
Brøchner, maaskee af personlige Grunde, foretrukket at oplyse Noget om Ordene.
% „Brøchner“ carries the ø slash here (verified at 600 dpi and at 8x). BOTH OCR
% witnesses drop it and read „Brochner“ -- tesseract loses every ø in the book,
% and the ABBYY layer fails on this one.
Oplysningen er „„at man ikke uden videre, fordi et Adjertiv kan prædiceres om et
% PRINTER'S DEFECT, p. 20, inside the quotation from Nielsen: the print reads
% „Adjertiv“ for „Adjectiv“. Read at 600 dpi and at 8x against the certain
% „Adjectiv“ two lines below on the same page: the fourth letter has the
% shoulder and the descender-less foot of Fraktur r, not the closed head of c.
% Transcribed as printed; both witnesses read „Adjertiv“ too. Recorded as a
% reading of the image, not as a claim about the sorts (JBIG2 --
% TRANSCRIPTION-PLAYBOOK §2).
Substantiv, derfor, omvendt, kan prædicere et til Substantivet svarende Adjectiv
om et til Adjectivet svarende Substantiv““, en Oplysning, som, hvad Sagen
angaaer, omtrent staaer paa lige Trin med Per Degns Oplysning om Imprimatur:
% „Imprimatur“ and the grammarians' terms „Nomen, Pronomen, Verbum,
% Participium, Conjugatio, Declinatio, Interjectio“ are FRAKTUR in the print,
% not antiqua; no \textit{} (verified at 600 dpi).
„„Alle de Ord, som nævnes kan, henføres til 8 Ting, som er Nomen, Pronomen,
Verbum, Participium, Conjugatio, Declinatio, Interjectio““. Denne Oplysning skal
% PRINTER'S DEFECT / QUOTE BALANCE, p. 20: the doubled inner quotation closes
% here with exactly two high marks („Interjectio““.“, counted at 8x), and the
% OUTER quotation -- the one opened at „Istedenfor above -- is never closed in
% the print. Transcribed as printed; the outer „ therefore has no partner and
% the fragment's quote counts do not balance. The two doubled pairs on this page
% are the only „„…““ in the book (RECON.md §3).
være min, men af Citatet af min Bog kan man see, at Oplysningen netop er Hr.
Prof. Nielsens. Jeg maa altsaa afstaae ham Æren af at have talt som Per Degn; men
jeg paaskjønner hans gode Villie, og skal til Gjengjæld om lidt være ham til
Tjeneste med et Citat fra det samme holbergske Stykke, til hvilket jeg kan tænke
mig, at Prof. N. hendrages ved en naturlig Sympathie. Men da jeg er kommen ind
paa Prof. N.'s Anmærkning, skal jeg tillade mig at tilføie en lille Bemærkning.
Prof. N. er bleven meget vred over, at jeg har talt om hans Forsøg i det
Burleske, og han taler noget gaadefuldt om „Anstændighed“ og „ærbare, hysteriske,
arrige gamle Skolejomfruer“, samt om, at der ikke i hans Exempler findes „urene
Bestanddele“, hvilke Ord han synes at tillægge mig som mine, siden han anfører
dem med Citationstegn, skjøndt jeg ingensteds har brugt dem. Professoren synes
% --- p. 21 ---
at være bleven lidt usikker paa sine æsthetiske Bestemmelser, og det kan jo let
hændes en Mand, der skal studere \textit{«practica, didactica, tactica,
homiletica, exegetica, ethica, rhetorica, oratoria, metaphysica, chiromantica,
necromantica, logica, talismanica, juridica, parasitica, politica, astronomica,
geometica, arithmetica»} og meget Mere, at Et og Andet kan gaae ham af Glemme.
% ANTIQUA with antiqua guillemets, the whole list, verified at 600 dpi.
% PRINTER'S DEFECT in it: the print reads „geometica“ for „geometrica“ (Holberg,
% Erasmus Montanus). It is broken across the line as „geo-“ / „metica“ -- with
% the ANTIQUA hyphen, not the Fraktur one -- and rejoined here; the r is absent
% at 8x. Transcribed as printed; both OCR witnesses read „geometica“.
Naar Æsthetiken atter kommer for Tour blandt Professorens mange Studier, vil det
være anbefaleligt for ham at opfriske „Begrebsudviklingen“ af det Burleske,
„Sagudviklingen“ vil han saa kunne finde i sine egne Skrifter.

Saa komme vi til Professorens Beskyldning imod mig for Mangel paa Agtelse for
Videnskaben. Denne Mangel godtgjøres ved et meget mærkeligt Beviis, hvis
Gyldighed atter støtter sig paa et endnu mærkeligere. Beviset, det uimodsigelige
Beviis, for min Mangel paa Agtelse for Videnskaben er det, at jeg ikke har brugt
mere end 6 Sider til at kriticere Grundideernes Logik, den epochegjørende
Grundideernes Logik! Jeg stod, og staaer endnu, i den Formening, at det var
fuldkommen tilstrækkeligt naar jeg kort og tydeligt paaviste, hvilket uopløseligt
Vilderede der i „Grundideernes“ Logik er i Grundideerne, altsaa netop i det, der
skulde holde Værket sammen, og naar jeg i et Par Exempler, tagne fra Punkter, der
have aldeles afgjørende Betydning for alt det Udkomne af Bogen, viste de
Tilsnigelser, ved hvilke Ordspilsdialektiken kom til sine Resultater. Men det er
efter Prof. N.'s Mening ingenlunde nok, og at det ikke er det, derfor skal saa
Beviset ligge i, at Grundideernes Logik er „ikke mindre end 60 Ark!“ Slaaende,
uomstødelige Beviis! Mærkelige videnskabelige Maalestok! Hvorfor angiver Hr.
Professoren ikke hellere Værkets Masse i Pund? Videnskabeligt var den Maalestok
ligesaa god, og praktisk vilde den være langt mere nyttig, især naar
Grundideernes Logik, efterat have ophørt at være tvungen Brugsgjenstand for de
yngre Studerende, samledes„ paa et eget Sted udenfor Literaturen“ med sine
% PRINTER'S DEFECT, p. 21: the opening „ is set tight against „samledes“, with
% the word space AFTER it instead of before it -- „samledes„ paa“ (verified at
% 600 dpi and at 8x). Transcribed as printed. Both OCR witnesses reproduce the
% anomaly; tesseract reads the mark „,,“, ABBYY the same.
forudgangne Søskende: den theologiske „propædeutiske Logik“ fra 1845 og
Prof. N.'s logiske Førstefødte, den stak\-%
% --- p. 22 ---
% „stakkels“ is broken by the compositor at the foot of p. 21 and rejoined here
% with \-% so that the transcription reflows.
kels lille Akephal, Systemet „uden Hoved og Hale“, den speculative Logik fra
1841, der „systematisk“ begynder med § 11 og ligesaa „systematisk“ ender midt i
en Sætning. Den \textit{petitio principii}, Prof. N. stadig gjør sig skyldig i
% ANTIQUA, no guillemets, verified at 600 dpi. „§ 11“ is set with a space and
% without a full stop after the section sign here; p. 26 has both „§ 35“ and
% „§. 11“ (RECON.md §8) -- the inconsistency is the book's.
med Hensyn til Grundideernes Logik, er den, at det overhovedet lønner sig Umagen
at beskæftige sig med dette Fragment af et logisk (anti-logisk?) Repertorium.
Havde Prof. N. ikke, fordybet som han er i sin Troeskritik og sin Videnskritik,
rent glemt, at der gives en tredie, meget nyttig Art Kritik, nemlig
\emph{Selv}kritik, saa vilde han maaskee have faaet et klarere Blik paa sine
% The compositor letterspaced only the first element: the print gives
% „S e l v kritik“, with „kritik“ set tight (verified at 600 dpi and at 8x).
% Encoded \emph{Selv}kritik, as p. 12's „lucus a n o n lucendo“ is encoded.
Præstationer. Prof. N. har velvilligt meddeelt mig Oplysninger om mit Talent,
maaskee vil en lignende Oplysning om hans eget Talent ikke være unyttig for ham.
Prof. N.'s videnskabelige Talent er et Talent for Specialundersøgelsen, hvor
Totalitetsanskuelsen og Grundtankerne ere ham \emph{givne}. Videnskabelig
Productivitet i høiere Betydning: Evnen til at frembringe de store, omformende
Tanker, mangler han ganske, og i samme Grad Evnen til at beherske og forme en
Heelhed. Forsøger han paa det Første, saa naaer han kun til Gjentagelser eller
Forquaklinger, og forsøger han paa det Sidste, saa opløser Heelheden sig strax
for ham i et ἄπειρον af ubehersket Detail. Derfor skulde Prof. N. aldrig have
% Greek, antiqua italic. Read off the image at 600 dpi and at 12x: the alpha
% carries the combined smooth breathing and acute, ἄ. The Fraktur model reads no
% Greek; the ABBYY layer reads the word but not its diacritics reliably.
taget fat paa et Arbeide som Grundideernes Logik, thi det \emph{maatte}
mislykkes for ham, fordi det fordrer begge de Evner, som Prof. N. saa
fuldstændig mangler. Hvorledes han savner Evne til at begrændse og forme et Stof,
det faaer man et slaaende Exempel paa i anden Deel af Grundideernes Logik. I
Fortalen til første Deel angiver Prof. N., at den indeholder „ikke fuldt to
Trediedele af den første Niendedeel af hele Værket“. Man skulde saa dog vente, at
det, der manglede i at udfylde dette „ikke fuldt“, laae nogenlunde overskuet for
Prof. N.'s Blik. Men idet han skal bearbeide det, svulmer det op for ham til en
saadan formløs Masse, at der, for at faae hiin Supplering tilveiebragt, foruden
hele anden Deel vil behøves endnu to Dele, hver paa 30 Ark, naar den samme
Maalestok skal beholdes. En uklar Følelse af sin Mangel paa Evne
% --- p. 23 ---
til Productivitet i den ovenfor angivne Betydning har Prof. N. havt, og den har
ført ham til, paa sine mange Udviklingstrin, eller rettere, under sine mange
Forvandlinger, altid at støtte sig til betydeligere Personligheder, først til
Hegel, saa til Martensen, saa til Kierkegaard, ligesom han nu, paa sit
% Hegel, Martensen, Kierkegaard, Grundtvig are all set in ordinary tight
% Fraktur here -- no letterspacing (measured at 600 dpi: median inter-letter gap
% 0.047--0.082 of the type height, against a page median of 0.068 and a
% letterspaced value of 0.17--0.27). Proper names are never antiqua in this book.
„overvidenskabelige“ Stadium, støtter sig til Grundtvig. I et saadant
Underordningsforhold vilde Prof. N. paa de enkelte Undersøgelsers Omraade kunne
udrette noget Godt og Nyttigt, saaledes som han har været nærmest ved at gjøre
det i sine Propædeutiske Forelæsninger fra 1860---61. Men naar Prof. N. vil være
% The mark after „61“ is a full stop, not a comma: measured on the image at
% 600 dpi it is 15 px wide and 20 px high and does not descend below the
% baseline, against a p. 23 comma of 18 x 29 px descending 13 px. The Fraktur
% OCR reads a comma here; the ABBYY layer reads the stop, and the image agrees
% with ABBYY. The same test settles „Logik“.“ and „Schema“.“ below.
selvstændig, og navnlig naar han vil være epochegjørende Reformator af
Videnskaben, saa gjør han Fiasco, som da han opfandt „den absolute
Uensartethed“ og da han skrev „Grundideernes Logik“.

Der er saa endnu kun een af Prof. N.'s Anker, jeg vil omtale for med det Samme
% NO PUNCTUATION after „omtale“ in the print, although the sense wants a comma.
% The Fraktur OCR reads a full stop there; what stands on the image is a speck
% 6 x 8 px, floating 6 px ABOVE the baseline, against 13--18 px for a p. 23 full
% stop and 14--21 x 24--33 px for a comma. It is dirt or a fleck of ink, not a
% sort, and is not transcribed. The ABBYY layer reads nothing there either.
at opfylde det Løfte, jeg nylig gav ham: Løftet om et brugbart Citat fra Erasmus
Montanus. Ved min upassende Behandling af de 60 Ark skal jeg ikke blot have viist
Mangel paa \emph{Agtelse} for Videnskaben, men ogsaa Mangel paa \emph{Sands} for
% LETTERSPACED, both confirmed twice over: spacing.py flags them (1.22 and 1.19)
% and the direct inter-letter measurement at 600 dpi gives 0.169 and 0.273 of
% the type height against a page median of 0.068. Their neighbours („Mangel“,
% „for“, „det rette“) measure 0.05--0.09 and are set tight.
det rette Videnskabelige. Egentlig skal nok Feilen ligge i, at jeg ikke har hævet
mig til Indsigt i „Forholdet mellem System og Schema“. Her skal jeg nu villig
indrømme, --- thi Noget bør man jo være saa billig at indrømme --- at jeg
virkelig ikke har kunnet naae til at tilegne mig Prof. N.'s Systematik. Men
Prof. N. har ogsaa øiensynligt udviklet det Systematiske til en Høide, der ikke
er kjendt i nogen tidligere Philosophie. Prof. N.'s philosophiske Skrifter ere
nemlig i en saa forbausende Grad gjennemtrængte af det Systematiske, at man kan
skjære dem over i større og mindre Stykker, og naar man saa sætter Stykkerne af
de forskjellige Bøger sammen efter Pagina og indbinder dem, saa er der strax et
System. Der manglede endnu kun, forat Prof. N. kunde naae det Systematiskes
„absolute Grændse“, at han indrettede sine Skrifter saaledes, at de, ligesom
visse af Middelalderens cabbalistiske Formler, kunde læses ligegodt bagfra
% --- p. 24 ---
\setcounter{footnote}{0}
og forfra. Forat nu Læseren ikke skal troe, at jeg gjør Løier, og forat Prof. N.
ikke atter skal bebreide mig, at jeg taler almindeligt i al Almindelighed om det
Almindelige, saa skal jeg nu give en „Sagudvikling“, saa „punktuel“, at Prof. N.
neppe kan ønske den punktuellere.
% Again a speck, not a stop, after „kan“ (5 x 4 px at 600 dpi); the Fraktur OCR
% reads „kan.“, the ABBYY layer nothing. Two more of the same on this page,
% after „til“ and after „Parti“; none is transcribed.

Prof. N. har i de sidste 3 Aar, i hvilke han har holdt philosophiske
Forelæsninger for de yngste Studerende, givet Forelæsningscursus i Boghandelen,
forskjellige for hvert Aar. I 1865---66 udkom et første og et andet Cursus, et
for hvert Halvaar. Det ene af disse Cursus blev dannet saaledes, at „Forel. over
% „ene“ is NOT letterspaced, though the raw inter-letter figure (0.163) looks
% high: „man“ nine lines below scores exactly the same 0.163 and is demonstrably
% tight (same width-to-height ratio as the tight „man“ in the footnote), so
% 0.16 is simply what three round Fraktur sorts give. spacing.py flags neither.
phil. Propædeutik fra 1860---61“ blev skaaret over ved pag. 177 og den
fraskaarne forreste Deel erstattet ved „Forel. over phil. Propæd. fra
1861---62“. Men nu havde denne sidste Bog kun 164 Pagina, og der skulde bruges
176. Der blev altsaa skrevet et Ark fuldt, der fandtes, passende at kunne give
% The comma after „fandtes“ is in the print (both witnesses), although the sense
% would rather have „der fandtes passende at kunne give Plads“. Kept as printed.
Plads til en Razzia mod Theologien; det blev skudt ind\footnote{% printed *)
  I dette indskudte Parti culminerer Prof. N.'s eiendommelige Vittighed: den
  Vittighed, der \emph{ikke} maa have noget \emph{Burlesk}. En elegant Prøve vil
  % LETTERSPACED, in the smaller footnote type: „ikke“ is 112 px wide against 82
  % for the tight „ikke“ nine lines below at the same type height, and
  % „Burlesk.“ measures 0.179 against a note median of about 0.065. The words
  % between them („maa have noget“) are tight -- the emphasis really is
  % discontinuous, on the two words only.
  man kunne finde S. 174, hvor Prof. N., vrængende ad den speculative Dogmatiks
  Skabelseslære, laver en Tabel med den „heelskabte“ Adam, den „næsten
  heelskabte og særdeles velskabte“ Eva, de „halvskabte men ikke vanskabte“
  Hverdagsmennesker, de „overhalvskabte“ Heroer og Genier, og de „halvoverskabte
  og just derfor ofte vanskabte“ dogmatiske Talenter. Smagfulde Vittigheder, der
  bydes de unge Studerende af den Lærer, der for dem skal repræsentere
  „Videnskabernes Videnskab!“}, to Blade omtrykte, og saa havde man et
„systematisk“ Værk. Hvad der var skaaret fra den større Bog, blev saa, efterat
det sidste Blad var omtrykt, til det andet Cursus, --- og var ligesaa
systematisk som før. Til de to Cursus føiedes der et Par korte Fortaler,
tilsammen paa 11 Sider, for det meste sammensatte af Brudstykker fra Prof. N.'s
philosophiske Propædeutik fra 1857 og fra Fortalen til Grundideernes Logik,
1ste Deel. Det maa vel være paa Grund af disse Fortaler, at
Prof. N.\footnote{% printed **)
  Jfr. Fortalen til det ene Cursus S. III}
% NO full stop closes this note in the print. What the Fraktur OCR takes for one
% measures 4 x 4 px at 600 dpi, against 12--13 px for the full stops in this
% same note („Talenter.“, „S.“, „Jfr.“). The ABBYY layer likewise reads none.
taler om hele Cursus'et som „disse \emph{Uddrag}, blandt
% --- p. 25 ---
hvilke ogsaa findes større og mindre Stykker af philos. Propædeutik,
Grundideernes Logik o. s. v.“ --- For Universitetsaaret 1866---67 udgaves et nyt
Cursus, der fik Titelen: „Propædeutik og Psychologie“. Det blev forfærdiget
saaledes. De propædeutiske Forelæsninger fra 1860---61 bleve atter deelte paa
% The Fraktur OCR puts a stray „-“ between „deelte“ and „paa“; it is a speck of
% the same kind as those on pp. 23 and 24 (the ABBYY layer reads nothing there)
% and is not transcribed. Per RECON.md §9 a lone „-“ from ABBYY may never be
% deleted without checking the image; the reverse case is checked the same way.
samme Sted som før, ved Side 177; men det Fraskaarne blev denne Gang erstattet
ved den philos. Propædeutik fra 1857. Her opstod saa den Vanskelighed, at denne
Bog havde 204 Sider, og der kunde kun bruges 176. Prof. N. vidste imidlertid,
paa en let og „systematisk“ Maade at hjælpe sig: han klippede Enden af Bogen, og
skjøndt det var de afsluttende og væsentligste Partier, der afklippedes, saa
blev „Systemet“ dog lige godt efterat to Blade ved Bogens Begyndelse og to ved
Sammenføiningen vare omtrykte. Som Supplement til dette nye Værk fik saa de
Studerende 2den Deel af Grundideernes Logik, og deri, at de maatte læse den,
% „2den“ is not letterspaced: the high raw figure comes from the side bearing of
% the figure 2, „Deel“ beside it is tight, and spacing.py does not flag it.
skjøndt de aldrig havde seet første Deel, var der naturligviis ogsaa „System“.
Men det Systematiske naaer dog først sin rette Høide i de to Cursus for
1867---68. Det ene: \emph{Logik og Psychologie}, er simplere i Sammensætningen.
% The two Cursus titles on this page are letterspaced, and the „Det ene:“ /
% „Det andet Cursus:“ that introduce them are not: „Logik“ and „Psychologie“
% measure 0.169 and 0.195, „Logik“ and „Propædeutik“ 0.220 and 0.200, against a
% page median of 0.064; „ene:“ scores 0.146, which is what a three-sort word
% plus a colon gives anyway (cf. the „ene“ / „man“ control on p. 24).
Side 1---262 er skaaret ud af Grundideernes Logik 1ste Deel, Side 279---482 af
de propæd. Forelæsninger fra 1860---61. Nyt er hvad der med særskilt Paginering
(S. 3---5) er sat foran Bogen, og Side 264---78, der betitles Psychologie, og
blandt Andet indeholder en kort og grundig Oversigt over det \emph{menneskelige}
% LETTERSPACED, and alone: „menneskelige“ measures 0.198 while „Sjæleliv,“ next
% to it measures 0.066. spacing.py flags it too (1.33). The irony is on the one
% word -- the whole of the HUMAN soul-life in two and a half pages.
Sjæleliv, paa 2½ Side, hvoraf dog et Stykke optages af en Afhandling om Sjælens
% Both witnesses read „21/2“; the printed form is the nut fraction 2½.
Udødelighed. Omtrykte ere kun de Blade, hvor Sammenføiningen er skeet.
Interessant er det, at det \emph{Selvsamme}, der i 1860 var \emph{Propædeutik},
i 1866 forestiller \emph{Psychologie}, i 1867 \emph{Biologie og Psychologie},
% The letterspacing here is DISCONTINUOUS and is set as measured: „Selv-“ 0.220,
% „samme,“ 0.185, „Propædeutik,“ 0.200, „Psychologie,“ 0.205, „Biologie“ 0.207,
% „Psychologie,“ 0.191 -- while „der“ 0.068, „var“ 0.067, „forestiller“ 0.056 and
% the year figures are all tight. „Selvsamme“ is broken across the line in the
% print and both halves are spaced.
--- men det er naturligviis ifølge en høiere „Systematik“! --- Det andet Cursus:
\emph{Logik og Propædeutik}, er mere combineret. Her er S. 3---80 skaaret fra de
propæd. Forelæsninger fra 1861---62, S. 82---177 fra dem fra 1860---61,
S. 177---204 fra den phil. Propædeutik fra 1857, og saa igjen S. 204---274 fra
Forelæsningerne fra 1860---61.
% --- p. 26 ---
Nyt er et „\emph{Grundrids af den formelle Logik}“ paa 2 (siger: to) Sider, der
% LETTERSPACED -- a third Cursus title spaced like the two on p. 25. Every
% measurable word of the printed line scores 0.15--0.18 against a page median of
% 0.060: „Grundrids“ 0.175, „den“ 0.17, „formelle“ 0.181, „Logik“ 0.151. The
% quotation marks are left outside the \emph{}.
er sat foran; omtrykte ere de 3 Blade, hvor Sammenføiningen er skeet. Man bliver
meget overrasket ved, midt i en Udvikling, der ikke er inddeelt i Paragrapher,
pludselig S. 184 at støde paa § 35 og S. 190 paa § 36, og saa igjen at see
Paragraphinddelingen ophøre. Overraskelsen hæves dog, naar man erindrer, at der
i forrige Aar var skaaret en Stump af den lille phil. Propædeutik, og det havde
jo været „usystematisk“, ikke at benytte den for de Paragraphers Skyld. Den
maatte jo ogsaa vække et kjærligt Minde hos Forfatteren om det lille System, der
begyndte med §. 11. --- Læseren vil nu, hvis han har Sands for det
% THE BOOK'S OWN INCONSISTENCY, kept: „§ 35“ and „§ 36“ above carry no full stop
% after the section sign, „§. 11“ here does. All three stand in one paragraph.
% (p. 22 has „§ 11“, without the stop -- see the pp. 13--22 batch.) The stop
% after the numeral „11“ is a separate, ordinary sentence stop: measured 16 x 18
% px at 600 dpi, the same as the stop on the section sign itself.
„Systematiske“, kunne udkaste Planen til Systemet i Prof. N.'s næste Cursus.
I Behold haves nemlig: af de propæd. Forelæsn. fra 1860---61: S. 1---82 og
S. 177---204, af dem fra 1861---62: S. 81---164, og af Grundideernes Logik 1ste
Bind: S. 263---Slutningen. Her behøve altsaa kun et Par smaa Mellemrum at
udfyldes, og naar de ligeligt anvendtes paa de „halvoverskabte og just derfor
ofte vanskabte“ dogmatiske Talenter, paa de opstyltede „Metaphysikmennesker“, og
paa de „philosophiske Dompapper“, saa kunde det jo blive meget „systematisk“.
Prof. N.'s Polemik vilde fortræffeligt egne sig til at sammenkitte Stykkerne til
et System; thi den er meget klæbrig. Til dette nye Cursus var det nu, at jeg
vilde anbefale Hr. Professoren det lovede Citat af Holberg. Hvad mener Hr.
Prof. N. om at sætte følgende Motto paa sit nye Cursus: „\emph{Naar denne Bog
% The Erasmus Montanus motto is NOT displayed and NOT indented: it runs straight
% on from „Cursus:“, whose line fills the measure, and the paragraph ends with
% „4de Scene.)“ in the middle of a line. It is marked by letterspacing alone,
% and the letterspacing is uniform over the whole quotation (0.13--0.27 against
% a page median of 0.060), the three antiqua words included. It fills five
% printed lines, not the seven RECON.md §5 estimates.
bindes ind i Tyrkisk Bind, saa kaldes den \textit{Donat}, men naar man binder
den ind i hvidt Pergament, saa kaldes den \textit{Grammatica} og declineres
ligesom \textit{ala}}“. (Erasmus Montanus 1ste Act., 4de Scene.)
% ANTIQUA: „Donat“, „Grammatica“ and „ala“ stand in upright antiqua against the
% Fraktur, and carry no guillemets. All three fall inside the letterspaced
% motto, so all three are spaced antiqua; the enclosing \emph{} carries that.
% The mark after „ala“ “ is a full stop (15 x 18 px), not the comma the Fraktur
% OCR reads. „Act.,“ carries both a stop and a comma, as RECON.md §6 records.

% Centred rule: a single hairline, 1.93 cm long (455 px at 600 dpi) and 9 px
% thick, with white space above and below, standing between the motto and the
% next paragraph. Text, not a footnote separator; there is no note on this page.
\begin{center}\rule{2cm}{0.4pt}\end{center}

Jeg er nu færdig med det, som man, hvis man vil tage Videnskab i „videre
Forstand“, kan kalde den „videnskabelige“ Side af Prof. N.'s Polemik; jeg gaaer
nu over
% --- p. 27 ---
til en anden Side af denne Polemik, og her maa jeg tale et noget andet Sprog.
Saalænge vi havde at gjøre med Prof. N.'s videnskabelige og literære Færden,
have vi ikke altid kunnet blive indenfor det Alvorliges Grændser: Gjenstanden
laae for langt inde paa Modsætningens Gebet. Her gjaldt det virkeligt:
\textit{difficile est satiram non scribere}. Men jeg kommer nu med Prof. N. ind
% ANTIQUA, upright, no guillemets -- as RECON.md §5 records. „non“ is NOT
% letterspaced here, unlike p. 12's „lucus a n o n lucendo“: its raw
% inter-letter figure is high only because n-o-n is two round gaps, and its
% width relative to „est“ on the same line (124 px against 96 px at a common
% x-height) is if anything narrower than the two words' natural advance ratio.
paa det Ethiskes Gebet, og her skal tales alvorligt. Det er om Prof. N.'s
Insinuation angaaende Bevæggrundene til min Optræden, at jeg nu skal tale et Ord.

Der er i de senere Aar oftere bleven ført Klage over den Form, i hvilken den
videnskabelige Polemik hos os hyppigt føres, og med Rette. Man har med Rette
klaget over den plumpe Tone, der viste, at vore Videnskabsmænd ikke altid tillige
vare dannede Mænd. Men Polemiken har i den senere Tid ikke sjelden faaet en uden
Sammenligning værre Charakteer, idet den er kommen ind paa at underskyde
Modstanderne forkastelige Motiver, og mod denne Art af Polemik, der især hyppig
er brugt mod Prof. N.'s Modstandere, skal der nedlægges en alvorlig Protest. Da
Cand. \emph{Larsen} i sin Tid skrev mod Prof. N., blev der i et Tidsskrift
anonymt insinueret, at Motivet havde været --- et Universitetsstipendium. Da
Dr. \emph{Heegaard} nylig skrev sit Skrift om Prof. N.'s Lære om Tro og Viden,
% LETTERSPACED, and Fraktur, not antiqua: 0.171 and 0.169 against a page median
% of 0.066. RECON.md §5 is right that no proper name in this book is set in
% antiqua; these two are spaced instead.
saa blev i en anonym Piece noget Lignende plumpt antydet. Og nu da jeg er
fremkommen med min Kritik, fremsætter Prof. N. i sin Piece en Insinuation, af et
andet Indhold end hine, men ikke mindre fornærmelig. Prof. N. vil, efter den
bekjendte Vending, \emph{ikke} sige, at det ikke er Kjærlighed til Videnskaben,
% Only the FIRST „ikke“ is spaced: 144 px wide against 105--107 px for every
% tight „ikke“ on this page, while the second, four words later, measures 106.
men et ulmende Had til den positive Christendom, der har ledet mit Angreb paa
ham; nei tværtimod! han vil finde det rigtigt, om der tages kraftigt til
Gjenmæle, hvis Sligt bliver sagt. Prof. N.'s vil \emph{kun} --- i andre Ord sige
% PRINTER'S DEFECT, p. 27: the print reads „Prof. N.'s vil kun“, where the sense
% plainly wants „Prof. N. vil kun“. Both OCR witnesses read the genitive 's, and
% so does the image. Transcribed as printed; not corrected.
% „kun“ here is spaced (128 px) and the „kun“ in the next sentence is not (108).
\emph{det Selvsamme}! Han vil kun sige, at det nok \emph{ikke} er af
% „det“ 113 px against 91--96 px for the six other instances of „det“ on the
% page; „Selvsamme!“ and the second „ikke“ likewise spaced (0.192; 141 px).
\emph{videnskabelige} Bevæggrunde, men af en utaalmodig hastende Iver forat
\emph{„komme Aabenba\-%
% --- p. 28 ---
% „Aabenbaringstroen“ is broken by the compositor at the foot of p. 27 and
% rejoined here with \-% so that the transcription reflows. The letterspacing
% and the quotation both run across the page break: „komme“ 0.185, „Aabenba-“
% 0.239 on p. 27, „ringstroen“ 0.159 and „tillivs“ “ 0.239 on p. 28.
\setcounter{footnote}{0}%
ringstroen tillivs“}, at jeg er optraadt\footnote{% printed *)
  Skulde den besynderlige Slutning, Prof. N. (S. 11) udleder af mine Ord,
  maaskee sigte hen i samme Retning? Jeg havde i mit Skrift udviklet, at den
  Bevidsthedsform, der var særegen for de positive Religioners Trin, lod en
  Egoismens Rest blive tilbage selv i det religiøse Kjærlighedsforhold. Dette
  bliver nu hos Prof. N. til at Evangeliet skal være et „Foster af skjult
  Egoisme“ og „Helten fra Gethsemane den største Egoist, der har levet paa
  Jorden“. Jeg har antaget Prof. N.'s Slutning for en stor Tankeløshed; jeg
  vilde nødig ansee den for noget Værre.}. Prof. N., der saa tidt klart har
skildret Fanatismens onde Instincter, betænker sig ikke paa, at vække alle disse
onde Instincter ved at henkaste en Insinuation, paa hvis Sandhed han ikke selv
kan troe! Lad os dog, naar vi ville bære Hædersnavnet af Videnskabsmænd, først og
fremmest mindes den Forpligtelse, det paalægger os: at handle som det sømmer sig
for Videnskabsmænd; men en Henkasten af den Art Insinuationer kan idetmindste
jeg ikke anerkjende for sømmelig.

Men Prof. N. fører jo et Beviis, „et Indiciebeviis“, for sin Sigtelse. Dersom det
virkelig havde været Videnskabens Ret, der havde ligget mig paa Sinde, saa skulde
jeg blot --- have forhandlet med Prof. N. om Problemerne i Grundideernes Logik!
Altsaa! naar jeg vilde have hævdet Videns og Videnskabens Ret ligeoverfor
Theologien, og ligeoverfor en „Troes-Theorie“, der i lige Maade forvansker Viden
% The hyphen of „Troes-Theorie“ falls at the end of a printed line, so it is
% both the compound's hyphen and the break. It is kept, because the second half
% is capitalised („Theorie“): a compositor breaking the one-word „Troestheorie“
% -- the form used on pp. 3--12 of this book -- would have set „theorie“.
og Tro, saa skulde jeg have fordybet mig i uendelige logiske Undersøgelser med
denne Theories Repræsentant, medens han uhindret udbredte Forvirring og krænkede
Videnskabens Ret; --- jeg skulde altsaa med eet Ord have hævdet Videnskabens Ret
ved \emph{ikke} at hævde den. Hele dette afsindige „Indiciebeviis“ faaer kun
% 136 px against 105--106 px for the tight „ikke“ elsewhere on the page.
Mening naar en Mellemsætning suppleres, som vel maa have foresvævet Prof.
Nielsen. Med den vilde Beviset lyde saaledes: hvis det havde været Prof.
Brøchners Hensigt at gjøre Videnskabens Ret gjældende, saa maatte han ---
\emph{eftersom Prof. N. er Eet med Videnskaben, er dens Incarnation} --- have
% LETTERSPACED over the whole parenthesis: „eftersom“ 0.160, „Prof.“ 0.177,
% „Videnskaben,“ 0.215, „dens“ 0.311, „Incarnation“ 0.228, against a page median
% of 0.065. RECON.md §5 records the same phrase.
fordybet sig i hine grundlogiske Undersøgelser, og ladet den personliggjorte
Videnskab
% --- p. 29 ---
ene tale Videnskabens Sag. Saaledes kom der en Art Mening ind i Beviset, men,
Prof. N. tilgive mig! jeg kan ikke skjønne rettere end at den kun kom ind ved
det, i hvilket der \emph{ikke} var Spor af Mening.
% 131 px against 105 px for the tight „ikke“ two lines above.

Med denne Indsigelse mod Prof. N.'s uberettigede Insinuation skilles jeg fra ham.
Da jeg skrev imod ham, vidste jeg forud, af hvad Art hans Svar vilde blive hvis
et saadant kom; jeg havde „Indicier“ nok. Paa Prof. N.'s Tiltale har han nu
faaet Gjensvar; og om han end, i Trangen til at beholde det sidste Ord, skulde
skrive en ny Piece, vil jeg være fritagen for at tage nogetsomhelst Hensyn til
den, medmindre den baade i Form og Indhold er helt forskjellig fra den sidste.

% Centred rule: a single hairline, 1.92 cm long (455 px at 600 dpi) and 7 px
% thick, with white space above and below. It is faint on the render but it is
% there -- a mechanical sweep of the leaf finds an isolated horizontal ink run
% of that length at y = 1883--1889, between the two paragraphs. Text, not a
% footnote separator; this page carries no note.
\begin{center}\rule{2cm}{0.4pt}\end{center}

Endnu kun et Par Ord til mine Læsere vil jeg tilføie om min personlige Stilling
til det Problem, jeg har behandlet. Det er næsten 30 Aar siden, da jeg endnu var
ganske ung, at Problemet om Troens Forhold til Viden første Gang viste sig for
mig i sin hele Betydning og blev en Gjenstand for min alvorlige og stadige
Overveielse. Jeg førtes snart hen til en Opfattelse, der ikke laae fjernt fra
den, jeg i mit Skrift har udtalt, og der kom et Øieblik, da det blev min Pligt,
aabent at vedkjende mig min Overbevisning. Den blev da udtalt uforbeholdent og
uden Sky, skjøndt denne Udtalelse maatte gribe dybt ind i alle mine personlige
Forhold, og lægge maaskee uoverstigelige Hindringer iveien for en
Fremtidsvirksomhed for mig. Det var i Slutningen af 1841; i 1843, netop i disse
Dage for 25 Aar siden, skrev jeg saa en lille Bog mod Martensen i Anledning af
hans Afhandling om Daaben. Fra den Tid af og indtil jeg for faa Maaneder siden
udgav mit Skrift, har jeg Intet skrevet, der angik det Problem, der uafbrudt har
beskæftiget mig; og i min nu allerede mangeaarige Universitetsvirksomhed har jeg
kun i een Forelæsning (i 1858---59) behandlet det fra den historiske Side. Jeg
% No full stop after the closing parenthesis: the sentence runs on. The Fraktur
% OCR supplies one; the image and the ABBYY layer do not.
havde altsaa intet „Hastværk“; jeg lod rolig min Overbevisning modnes hos
% --- p. 30 ---
mig, indtil det blev mig sikkert, at den ikke blot var videnskabelig holdbar,
men at den kunde bære et Personlighedens Liv. Og jeg vilde ikke vove at paatage
mig det Ansvar, ved en utidig Meddelelse kun at nedbryde hvor jeg Intet kunde
opbygge. Men der kom atter et Øieblik, da jeg ansaae det for min Pligt, ikke at
tie. Det var da det videnskabeligt Værdiløse og det personligt Uvederhæftige
anmassende trængte sig frem, udspredende Forvirring, og da det af tankeløs
Ukyndighed markskrigerisk udraabtes som den eneste Sandhed. Jeg gjorde saa
Problemet til Gjenstand for den Universitetsforelæsning, der nu foreligger
offentliggjort. I hvilken Aand og Tone det er behandlet, derom har Læseren nu
Leilighed til selv at dømme; og jeg tør sikkert stole paa, at Ingen skal kunne
paavise Spor af Had eller af Mangel paa Ærefrygt for det, der er det Helligste
for saa Mange; jeg tør stole derpaa; thi fra alvorlige og ædle Mænd, inderligt
knyttede til det positivt Religiøse, har jeg modtaget Udtalelser, fulde af
Velvillie og Anerkjendelse. Og jeg tør roligt rette det Spørgsmaal til alle dem,
der i de mange Aar have været mine Tilhørere, ogsaa til dem, der meest
lidenskabeligt have knyttet sig til en Overbevisning, der er min modsat: om jeg
nogensinde som Universitetslærer har udtalt et Ord, der kunde saare deres
Følelser, om jeg ikke tværtimod bestandig har talt med Ærbødighed om
Christendommen, om jeg ikke redeligt har stræbt at vise dens overlegne Betydning
i Forhold til de Aandsphænomener, der ved dens Fremkomst fremgik af en lignende
Aandens Stræben. Fordi jeg ikke søger Sandheden i det, der nu i Aarhundreder har
baaret saa mange Menneskeslægters Liv, men i det, der \emph{evigt} bærer
% The one letterspaced word on the page: 0.173 against a page median of 0.060,
% and spacing.py flags it too (1.31). Its neighbours are tight.
Menneskeaandens Liv, og knytter det sammen i Enhed med dets guddommelige
Udspring, skulde jeg derfor ikke forstaae at agte de Former, gjennem hvilke
Menneskeaanden stræbte hen mod denne Enhed, og navnlig den reneste og aandigste
af disse Former? Tilvisse! det kan Ingen mene, og selv min Modstander kan ikke
troe det; thi vi have tidt nok udvexlet vore Tanker om denne Sag.

Min Opgave ligger klar for mig, og uforstyrret af An\-%
% --- p. 31 ---
% „Angreb“ is broken by the compositor at the foot of p. 30 and rejoined here
% with \-% so that the transcription reflows.
greb skal jeg stræbe at løse den. Den er: ad Videns Vei at finde det, der
bringer Enhed og Forsoning i Menneskeaandens Liv, Forsoning mellem Mennesket og
dets Udspring. Paa Løsningen af denne Opgave skal jeg arbeide nu som før, uden
Lidenskab og uden Had, glædende mig ved de alvorligt og redeligt Stræbendes
Sympathie, med Sandheden til min eneste Norm, altid paavirkelig ved klare
Grunde, men uskræmmet ved Fanatismens Larmen og med rolig Overlegenhed overfor
det Lave.

% Ornamental tailpiece closing the book: a centred printer's flower, 4.06 cm
% wide (959 px at 600 dpi) and 43 px tall -- a swelled rule tapering to a point
% at each end, with a chain of about a dozen lozenge beads through the middle.
% Set here like the two hairlines on pp. 26 and 29, since the repo has no
% ornament stock. Nothing follows it on the leaf; there is no colophon.
\begin{center}\rule{2cm}{0.4pt}\end{center}

% =====================================================================
%  RETTELSE  (printed on a leaf at the back, PDF 37)
%  Fraktur. The heading is letterspaced and centred; the line below it is
%  centred and not spaced. Transcribed as printed; NOT applied to p. 11.
% =====================================================================
\bigskip
\begin{center}
  \emph{Rettelse.}

  \smallskip
  Side 11, Lin. 2: videnskabelige, læs: subjective.
\end{center}

\end{document}
